\documentclass[10pt]{article}

\usepackage[a4paper,margin=2.2cm]{geometry}
\usepackage[T1]{fontenc}
\usepackage[utf8]{inputenc}
\usepackage{lmodern}
\usepackage{microtype}
\usepackage{amsmath,amssymb,mathtools}
\usepackage{booktabs,tabularx,array,multirow}
\usepackage{enumitem}
\usepackage{xcolor}
\usepackage{graphicx}
\usepackage{tikz}
\usetikzlibrary{arrows.meta,positioning,fit,shapes.geometric}
\usepackage{listings}
\usepackage{hyperref}
\usepackage[nameinlink,noabbrev]{cleveref}
\usepackage{caption}
\usepackage{subcaption}
\usepackage{url}
\usepackage{balance}

\definecolor{acmblue}{RGB}{37,80,130}
\definecolor{acmgray}{RGB}{245,247,249}
\hypersetup{
  colorlinks=true,
  linkcolor=acmblue,
  citecolor=acmblue,
  urlcolor=acmblue,
  pdftitle={Agentic Configuration Management},
  pdfauthor={AUTHOR NAME}
}

\lstdefinestyle{acmcode}{
  basicstyle=\ttfamily\small,
  backgroundcolor=\color{acmgray},
  frame=single,
  rulecolor=\color{black!15},
  breaklines=true,
  columns=fullflexible,
  showstringspaces=false
}

\newcommand{\acmcore}{\textsc{ACM Core}}
\newcommand{\framework}{framework-independent}
\newcommand{\todo}[1]{\textcolor{red}{[TODO: #1]}}
\newcommand{\system}{\mathcal{S}}
\newcommand{\configgraph}{G_C}
\newcommand{\runtimegraph}{G_R}
\newcommand{\assurancegraph}{G_A}
\newcommand{\evolutiongraph}{G_E}

\title{\textbf{Agentic Configuration Management (ACM):\\A Reference Configuration Model for Governed Agentic Systems}}

\author{
  AUTHOR NAME\\
  AFFILIATION\\
  \texttt{EMAIL ADDRESS}
}

\date{Draft --- July 2026}

\begin{document}
\maketitle

\begin{abstract}

Agentic systems are rapidly evolving from isolated LLM applications to complex software composed of interacting agents, tools, prompts, models, and execution workflows. While existing LLMOps and AgentOps platforms provide valuable support for orchestration and observability, they do not offer a framework-independent configuration model capable of describing, validating, and governing the system as a coherent whole. As a result, configuration management remains fragmented across frameworks, making reproducibility, impact analysis, and long-term maintenance increasingly difficult.

This paper introduces \textbf{Agentic Configuration Management (ACM)}, a reference configuration model designed for governed agentic systems. ACM extends established Software Configuration Management principles to the agentic domain through immutable configuration items, explicit baselines, lifecycle management, dependency-aware state propagation, and runtime provenance. Rather than replacing existing orchestration frameworks, ACM provides a common representation that enables their configurations to be normalized and managed consistently.

We describe a complete reference implementation in Python together with adapters for LangGraph and CrewAI. The implementation includes a deterministic validation engine, a replayable event model, and a preservation-based extraction process that reconstructs native workflows into the ACM representation. The evaluation combines normative scenarios with automated tests to assess both the consistency of the model and the fidelity of the extraction process across frameworks.

The results show that ACM can preserve the structural information required for configuration governance while making explicit the information that is approximated or unavailable in the underlying framework. These results provide initial evidence that a unified configuration model can support reproducibility, auditability, and governance across heterogeneous agentic ecosystems without being tied to a specific orchestration framework.

\end{abstract}

\section{Introduction}

Recent advances in large language models (LLMs) have accelerated the emergence of \emph{agentic systems}, in which autonomous software agents cooperate to accomplish complex tasks through reasoning, tool invocation, delegation, and dynamic workflow execution. Unlike traditional software applications, these systems are no longer defined solely by source code. Their behaviour depends on a heterogeneous set of artefacts, including prompts, models, tools, policies, workflows, evaluation profiles, and runtime-generated entities, whose interactions evolve throughout the system lifecycle.

This evolution challenges the assumptions underlying traditional Software Configuration Management (SCM). Classical SCM standards were designed for software systems whose configuration is primarily composed of relatively static artefacts, such as source files, libraries, executables, and deployment descriptors. Although these principles remain essential, they do not directly address the mutable and heterogeneous nature of modern agentic systems, where configuration may evolve during execution, new agents may be instantiated dynamically, and behavioural changes can arise without modifications to the application source code itself.

In parallel, the rapid development of LLMOps, AgentOps, and agent orchestration frameworks has significantly improved the engineering of agentic applications. LLMOps platforms provide observability, evaluation, prompt management, and experimentation capabilities, while orchestration frameworks such as LangGraph and CrewAI facilitate the construction and execution of multi-agent workflows. However, these approaches primarily focus on operational concerns or framework-specific execution models. They do not provide a framework-independent representation capable of describing the complete configuration of an agentic system, managing its lifecycle, tracing its evolution, or supporting reproducible governance across heterogeneous ecosystems.

This fragmentation has practical consequences. The absence of a common configuration model complicates configuration auditing, impact analysis, reproducibility, and long-term maintenance. It also limits interoperability between frameworks, since each ecosystem relies on its own internal abstractions to represent agents, workflows, prompts, tools, and execution semantics. As agentic systems become increasingly dynamic and are progressively deployed in production environments, configuration governance itself becomes a first-class engineering concern.

This paper argues that agentic systems require a dedicated configuration abstraction extending the principles of Software Configuration Management while remaining independent of any execution framework. Rather than introducing yet another orchestration framework, we propose \emph{Agentic Configuration Management (ACM)}, a reference model for representing, governing, and auditing agentic systems throughout their lifecycle. ACM introduces a normalized representation of agentic configuration items, immutable revisions, explicit baselines, provenance-aware relationships, lifecycle semantics, runtime state reconstruction, and assurance-oriented governance. These concepts provide a common language for describing agentic systems independently of the technologies used to implement or execute them.

To demonstrate the feasibility of the proposed model, we present a reference implementation in Python together with extraction adapters for LangGraph and CrewAI. The implementation operationalizes the ACM specification through deterministic validation, lifecycle management, event replay, state propagation, and framework normalization. Rather than evaluating the implementation as a software product, the experiments assess the expressiveness of the model, the fidelity of framework projection, and the consistency of its governance semantics across representative agentic workflows.

\begin{figure*}[t]
	\centering
	\includegraphics[width=\textwidth]{figures/ACM_architecture_principles.png}
	\caption{Positioning of ACM.ACM acts as a framework-independent configuration layer that connects agent frameworks (orchestration), LLMOps platforms (observability and evaluation) and deployment environments (execution).}
	\label{fig:ACM_positioning}
\end{figure*}

The contributions of this paper are fourfold:

\begin{itemize}
    \item We identify the limitations of existing Software Configuration Management, LLMOps, AgentOps, and agent orchestration approaches with respect to the governance of modern agentic systems.
    
    \item We introduce Agentic Configuration Management (ACM), a framework-independent reference model defining configuration artefacts, relationships, lifecycle semantics, provenance, runtime representation, and assurance mechanisms for governed agentic systems.
    
    \item We describe a reference implementation that operationalizes the ACM model through deterministic validation, state propagation, baseline management, and framework projection while preserving the conceptual independence of the model from its implementation.
    
    \item We evaluate the proposed model through normative scenarios and cross-framework projection experiments using LangGraph and CrewAI, providing initial evidence of its expressiveness, consistency, and portability across distinct orchestration paradigms.
\end{itemize}

The remainder of this paper is organized as follows. Section~\ref{sec:related_work} reviews the foundations of Software Configuration Management together with recent advances in AI governance, LLMOps, AgentOps, and agent orchestration frameworks. Section~\ref{sec:design_principles} derives the design principles that motivate ACM. Sections~\ref{sec:reference_model} and~\ref{sec:governance_semantics} present the reference model and its governance semantics. Section~\ref{sec:operationalization} describes the reference implementation, while Section~\ref{sec:evaluation} reports the experimental evaluation. Finally, Sections~\ref{sec:discussion}, \ref{sec:threats}, \ref{sec:future_work}, and \ref{sec:conclusion} discuss the implications, limitations, future directions, and conclusions of this work.


\section{Related Work}
\label{sec:related_work}

\subsection{Foundations of Software Configuration Management}
\label{sec:scm_foundations}

Software Configuration Management (SCM) emerged as a response to the growing complexity of software systems and the need to control the evolution of their constituent artefacts throughout the software lifecycle. As software projects increased in size and involved multiple developers, uncontrolled modifications rapidly became a major source of integration failures, reproducibility issues, and maintenance costs. SCM was therefore established to ensure that software configurations could be uniquely identified, consistently reproduced, systematically modified, and independently audited.

The discipline has progressively evolved from simple version control practices into a comprehensive configuration governance process. International standards such as IEEE~828 and ISO/IEC/IEEE~12207 define configuration management as a set of coordinated activities supporting the identification of configuration items, the control of changes, the recording of configuration status, and the verification that software products conform to their approved specifications. These activities have remained remarkably stable despite the evolution of software engineering practices, reflecting their fundamental role in maintaining the integrity and traceability of complex systems.

A central concept of SCM is the \emph{Configuration Item} (CI), defined as an artefact placed under configuration control because of its importance to the system. Configuration items may include source code, libraries, executable binaries, documentation, configuration files, interface specifications, test artefacts, or deployment descriptors. Each configuration item is uniquely identified, versioned, and managed independently, allowing changes to be introduced in a controlled and traceable manner.

SCM also introduces the notion of a \emph{baseline}, which represents an approved and immutable snapshot of a collection of configuration items at a particular point in time. Baselines establish stable references from which software systems can be built, deployed, reproduced, or audited. Once released, a baseline is modified only through a formally controlled change process, ensuring that every evolution remains attributable to explicit configuration decisions.

Traditional SCM is commonly organized around four fundamental activities, illustrated in Figure~\ref{fig:scm_lifecycle}. Configuration identification establishes the set of managed artefacts together with their unique identifiers and relationships. Change control governs the introduction, evaluation, approval, and implementation of modifications. Configuration status accounting records the history and current status of configuration items and baselines, providing complete traceability of their evolution. Finally, configuration audits verify that approved configurations are internally consistent and faithfully correspond to their documented specifications.

\begin{figure*}[t]
	\centering
	\includegraphics[width=\textwidth]{figures/SCM_configuration_management_liefcycle_baseline.png}
	\caption{Caption SCM to complete.}
	\label{fig:SCM_framework}
\end{figure*}

These principles have demonstrated their effectiveness across several decades of software engineering practice and remain essential for ensuring reproducibility, maintainability, and governance in large-scale software projects. Importantly, they define a conceptual framework rather than a particular implementation technology. Consequently, they provide a solid theoretical foundation upon which configuration governance models for emerging software paradigms can be constructed.

\subsection{AI Governance}
\label{sec:ai_governance}

The rapid adoption of Artificial Intelligence (AI) across critical industrial, governmental, and societal domains has shifted the focus of software engineering beyond functional correctness toward broader governance concerns. Modern AI systems are increasingly expected to satisfy requirements related to transparency, accountability, traceability, reproducibility, and regulatory compliance throughout their operational lifecycle. These requirements have progressively established AI Governance as a complementary discipline whose objective is to ensure that AI systems remain trustworthy, controllable, and auditable from development to deployment.

Unlike traditional software governance, AI Governance must address the inherent variability of data-driven systems. Model behaviour depends not only on source code but also on training datasets, learned parameters, prompts, external tools, runtime context, and human interactions. Consequently, understanding why an AI system produced a particular result often requires reconstructing the complete decision context rather than inspecting the implementation alone. Governance therefore extends beyond software artefacts to encompass the entire chain of evidence supporting the behaviour of an AI system.

Several governance frameworks and regulatory initiatives have converged toward a common set of principles. Although they differ in scope and terminology, they consistently emphasize transparency, human oversight, accountability, risk management, documentation, and lifecycle traceability. These principles are reflected in recent international standards as well as emerging regulatory frameworks such as the European Union AI Act, highlighting the growing importance of maintaining reliable evidence regarding the development, deployment, and operation of AI systems.

A fundamental concept introduced by AI Governance is \emph{provenance}. Provenance describes the origin of an artefact together with the sequence of transformations that produced it. Closely related notions include \emph{lineage}, which captures dependency relationships among artefacts across successive transformations, and \emph{traceability}, which enables the reconstruction of decisions and modifications throughout the lifecycle of a system. Together, these concepts provide the evidence required to justify the current state of an AI system and to reproduce its behaviour under comparable conditions.

AI Governance also introduces the notion of continuous accountability. Rather than validating a system only at deployment time, governance activities accompany the entire operational lifecycle through monitoring, documentation, evaluation, incident analysis, and periodic reassessment. As AI systems evolve through successive revisions, governance mechanisms ensure that modifications remain attributable, documented, and consistent with previously established requirements.

These governance principles naturally complement the objectives of Software Configuration Management. While SCM provides mechanisms for identifying, controlling, and auditing software artefacts, AI Governance extends the scope of governance toward behavioural evidence, operational accountability, and lifecycle oversight. Together, they establish the conceptual foundations required for managing increasingly complex AI-based systems.

Nevertheless, current AI Governance frameworks generally remain technology-agnostic and operate at the organizational or regulatory level. They specify \emph{what} should be governed but rarely define \emph{how} the internal configuration of an AI system should be represented, versioned, or governed throughout its evolution. In particular, they provide little guidance regarding heterogeneous configuration artefacts such as prompts, agent definitions, workflow topologies, runtime-generated entities, or cross-component dependencies. This observation motivates the need for a configuration-centric reference model capable of operationalizing governance principles for modern agentic systems.

\subsection{LLMOps and AgentOps}
\label{sec:llmops_agentops}

The emergence of Large Language Models (LLMs) has profoundly transformed the software lifecycle. Unlike traditional machine learning systems, LLM-based applications are increasingly assembled from prompts, external tools, retrieval mechanisms, memory components, and autonomous reasoning workflows whose behaviour depends on runtime interactions rather than solely on deterministic program logic. As these applications evolved from isolated chat interfaces to production-grade systems, new engineering practices collectively referred to as \emph{LLMOps} have emerged to support their development, deployment, monitoring, and continuous improvement.

LLMOps extends the principles of MLOps by addressing the operational characteristics specific to foundation models. Beyond model deployment, LLMOps encompasses prompt management, execution tracing, experimentation, evaluation pipelines, cost monitoring, latency analysis, and continuous quality assessment. These practices aim to improve the reliability and maintainability of LLM-based applications while enabling iterative development through systematic observation of runtime behaviour.

As multi-agent systems have become increasingly widespread, the scope of operational concerns has expanded beyond individual model invocations. This evolution has led to the emergence of \emph{AgentOps}, which focuses on the engineering and operation of autonomous agentic systems. AgentOps extends LLMOps by considering interactions among multiple agents, orchestration workflows, delegated tasks, tool invocations, memory updates, and long-running executions. Consequently, observability must no longer capture isolated LLM calls but the behaviour of an entire distributed reasoning process.

Several concepts have progressively become fundamental within both LLMOps and AgentOps.

\paragraph{Configuration.}
Configuration describes the set of artefacts that determine the behaviour of an application independently of its source code. In contemporary LLM systems, configuration may include prompts, model parameters, retrieval settings, tools, evaluation datasets, execution policies, and deployment parameters.

\paragraph{Tracing.}
Tracing records the sequence of operations executed during a request. Modern tracing systems represent executions as hierarchical traces composed of spans corresponding to model invocations, tool calls, retrieval operations, or intermediate reasoning steps. Tracing enables debugging, performance analysis, and execution replay.

\paragraph{Observability.}
Observability refers to the capability of understanding the internal behaviour of a system through externally observable information. In LLM applications, observability combines traces, logs, metrics, costs, latencies, token consumption, and execution metadata to support operational diagnosis and optimization.

\paragraph{Provenance and Lineage.}
Although often used interchangeably, provenance and lineage capture complementary aspects of system evolution. Provenance identifies the origin of an artefact together with the transformations that produced it, whereas lineage describes dependency relationships across successive artefacts and executions. Both concepts are essential for reproducibility and accountability.

\paragraph{Evaluation.}
Evaluation has become a continuous process rather than a one-time validation activity. Current LLMOps platforms support automatic benchmark execution, human annotation, LLM-as-a-judge methodologies, regression testing, dataset management, and comparative experimentation across successive prompt or model revisions.

\paragraph{Deployment State.}
Operational platforms also maintain deployment metadata describing which prompt, model, or configuration version is currently active in a given environment. Such metadata supports rollback strategies, staged deployments, and production monitoring.

A rich ecosystem of platforms has emerged to operationalize these concepts. Commercial and open-source solutions such as LangSmith, LangFuse, Helicone, Braintrust, Weave, Phoenix, and PromptLayer provide comprehensive capabilities for tracing, observability, prompt versioning, experimentation, evaluation, and performance monitoring. Although their implementation strategies differ, they share a common objective: collecting operational evidence that facilitates debugging, optimization, and continuous improvement of LLM-based applications.

Despite their maturity, these platforms primarily operate at the execution level. Their central abstractions are traces, spans, prompts, evaluations, sessions, and runtime metrics. Prompt management capabilities have significantly evolved in recent years, introducing prompt versioning, deployment labels, rollback mechanisms, A/B experimentation, and evaluation datasets. Nevertheless, these artefacts are generally managed independently rather than as components of a unified system configuration.

Similarly, AgentOps platforms extend observability to multi-agent executions by reconstructing interaction graphs, monitoring delegated tasks, and recording execution histories. However, their representations remain tightly coupled to the execution model of the underlying orchestration framework. They describe \emph{how} a particular execution unfolded, but not necessarily the complete configuration that made this execution possible.

Consequently, current LLMOps and AgentOps solutions provide excellent operational visibility but do not define a framework-independent configuration model capable of representing heterogeneous configuration artefacts, managing immutable baselines, governing configuration evolution, or expressing dependency semantics independently of runtime execution. Their primary objective is to observe systems in operation rather than to provide a normalized representation suitable for long-term configuration governance.

\subsection{Agent Frameworks}
\label{sec:agent_frameworks}

The rapid adoption of agentic AI has led to the emergence of numerous orchestration frameworks that simplify the construction of systems composed of autonomous or semi-autonomous agents. Although these frameworks differ significantly in their programming models and execution semantics, they share a common objective: coordinating interactions among language models, external tools, memory components, and execution workflows to accomplish complex tasks.

Current agent frameworks generally provide abstractions for defining agents, orchestrating their interactions, invoking external tools, and managing execution state. Depending on the framework, these abstractions may be represented as explicit execution graphs, task-oriented workflows, hierarchical organizations of agents, or dynamic delegation mechanisms. Their primary concern is therefore the operational execution of agentic systems rather than the long-term management of their configuration.

Among open-source frameworks, LangGraph has become one of the most widely adopted platforms for developing stateful LLM applications. It models applications as explicit directed graphs whose nodes represent computational units and whose edges define possible execution transitions. Persistent shared state, checkpointing, human-in-the-loop interactions, and deterministic workflow execution make LangGraph particularly suitable for long-running agentic applications. Because its execution model is explicitly graph-based, the relationships among workflow nodes are directly represented within the framework.

CrewAI adopts a substantially different design philosophy. Instead of exposing an execution graph as the primary abstraction, CrewAI organizes applications around agents, roles, tasks, crews, and flows. Agents are assigned specialized responsibilities and collaborate through task delegation and orchestration mechanisms. This higher-level abstraction emphasizes organizational structure and collaborative behaviour rather than explicit graph topology, making CrewAI representative of role-oriented multi-agent systems.

Several additional frameworks have recently emerged to address complementary aspects of agent orchestration. The OpenAI Agents SDK introduces a lightweight programming model centered on agents, tools, guardrails, tracing, and handoffs, enabling dynamic delegation between specialized agents during execution. Microsoft Agent Framework integrates orchestration, state management, middleware, telemetry, and enterprise deployment capabilities within a unified platform. Google Agent Development Kit (ADK) provides a modular environment for constructing, evaluating, and deploying multi-agent systems, with particular emphasis on the Google AI ecosystem. Although these frameworks differ in their abstractions and runtime architectures, they all focus primarily on the execution of agentic applications rather than on the governance of their configuration.

Table~\ref{tab:framework_comparison} summarizes the principal characteristics of representative agent frameworks.

\begin{table*}[t]
\centering
\caption{Representative agent frameworks and their primary capabilities.}
\label{tab:framework_comparison}
\small
\begin{tabular}{p{2.6cm}p{2.6cm}cccccc}
\toprule
\textbf{Framework} &
\textbf{Primary abstraction} &
\textbf{Execution orchestration} &
\textbf{Tool integration} &
\textbf{State management} &
\textbf{Tracing} &
\textbf{Configuration versioning} &
\textbf{Configuration governance} \\
\midrule
LangGraph &
Explicit execution graph &
\checkmark &
\checkmark &
\checkmark &
Partial &
Limited &
No \\
CrewAI &
Agents, tasks, crews &
\checkmark &
\checkmark &
Partial &
Partial &
Limited &
No \\
OpenAI Agents SDK &
Agents and handoffs &
\checkmark &
\checkmark &
Partial &
\checkmark &
Limited &
No \\
Microsoft Agent Framework &
Workflows and agents &
\checkmark &
\checkmark &
\checkmark &
\checkmark &
Limited &
No \\
Google ADK &
Multi-agent workflows &
\checkmark &
\checkmark &
\checkmark &
Partial &
Limited &
No \\
\bottomrule
\end{tabular}
\end{table*}

Despite their architectural diversity, these frameworks exhibit several common characteristics. First, configuration information is typically distributed across multiple implementation artefacts, including source code, prompts, workflow definitions, deployment parameters, and framework-specific metadata. Second, configuration objects are generally managed according to framework-specific internal representations, limiting interoperability across ecosystems. Finally, although many frameworks provide execution traces and runtime telemetry, they rarely maintain an explicit, immutable representation of the complete system configuration throughout its lifecycle.

Consequently, existing agent frameworks provide sophisticated execution environments but stop short of defining a framework-independent configuration model capable of representing heterogeneous configuration artefacts, tracking their evolution, expressing dependency relationships, or supporting long-term configuration governance. These observations motivate the comparative analysis presented in the next section.

\subsection{Gap Analysis}
\label{sec:gap_analysis}

The previous sections have reviewed four complementary research and engineering domains that contribute to the governance of modern AI systems: Software Configuration Management (SCM), AI Governance, LLMOps/AgentOps, and agent orchestration frameworks. Although each addresses important aspects of the lifecycle of intelligent systems, their objectives, abstractions, and operational scope remain fundamentally different.

Software Configuration Management establishes rigorous principles for identifying, versioning, auditing, and controlling software artefacts. These mechanisms have demonstrated their effectiveness for conventional software systems but were designed for relatively static artefacts such as source code, documentation, and build configurations. They do not explicitly address the heterogeneous and highly dynamic artefacts that characterize contemporary agentic systems, including prompts, agent definitions, runtime-generated entities, or autonomous workflow evolution.

AI Governance approaches extend the scope of governance toward organizational accountability, transparency, risk management, regulatory compliance, and trustworthy AI. Their emphasis lies on defining governance principles and organizational processes rather than specifying how the internal configuration of an agentic system should be represented or managed. Consequently, these frameworks establish governance objectives but generally remain agnostic regarding the underlying configuration model.

LLMOps and AgentOps platforms focus primarily on operational excellence. Their major contributions include execution tracing, prompt management, evaluation pipelines, deployment monitoring, and observability. These capabilities substantially improve the operation of LLM-based applications but remain centered on execution evidence rather than on the governance of a complete system configuration. Configuration artefacts are generally managed independently and rarely represented as components of a coherent configuration model.

Finally, agent orchestration frameworks provide increasingly sophisticated abstractions for constructing and executing autonomous systems. Their primary objective is to define execution semantics, workflow orchestration, communication mechanisms, and runtime coordination among agents. Although several frameworks expose internal workflow representations, these representations remain framework-specific and are not intended to serve as interoperable configuration models spanning multiple execution environments.

Taken individually, each of these domains addresses an essential aspect of agentic systems. Taken together, however, they reveal the absence of a framework-independent reference model capable of representing the complete configuration of an agentic system while simultaneously supporting lifecycle management, provenance, dependency semantics, runtime evolution, and governance policies. Table~\ref{tab:gap_analysis} summarizes this observation.

\begin{table*}[t]
\centering
\caption{Comparison of existing governance approaches and the capabilities addressed by ACM.}
\label{tab:gap_analysis}
\small
\renewcommand{\arraystretch}{1.15}
\begin{tabular}{p{4.8cm}ccccc}
\toprule
\textbf{Capability} &
\textbf{SCM} &
\textbf{AI Gov.} &
\textbf{LLMOps / AgentOps} &
\textbf{Agent Frameworks} &
\textbf{ACM} \\
\midrule

\multicolumn{6}{l}{\textit{Configuration Representation}} \\
\midrule
Framework-independent configuration model & -- & -- & -- & -- & \checkmark \\
Configuration graph representation         & -- & -- & -- & Partial & \checkmark \\
Cross-artifact versioning                  & Partial & -- & Partial & -- & \checkmark \\
Dependency modelling                       & Partial & -- & Partial & Partial & \checkmark \\
Immutable baselines                        & \checkmark & -- & -- & -- & \checkmark \\

\addlinespace

\multicolumn{6}{l}{\textit{Governance}} \\
\midrule
Lifecycle management                       & \checkmark & \checkmark & Partial & Partial & \checkmark \\
Configuration policies                     & Partial & \checkmark & Partial & Partial & \checkmark \\
Configuration assurance                    & -- & Partial & Partial & -- & \checkmark \\
Configuration impact analysis              & Partial & Partial & -- & -- & \checkmark \\

\addlinespace

\multicolumn{6}{l}{\textit{Runtime Management}} \\
\midrule
Execution observability                    & -- & Partial & \checkmark & \checkmark & \checkmark \\
Runtime configuration evolution            & -- & Partial & Partial & Partial & \checkmark \\
Dynamic agent governance                   & -- & -- & Partial & Partial & \checkmark \\

\addlinespace

\multicolumn{6}{l}{\textit{Traceability}} \\
\midrule
Configuration provenance                   & Partial & \checkmark & Partial & Partial & \checkmark \\
Configuration lineage                      & Partial & Partial & Partial & Partial & \checkmark \\
Auditability and reproducibility           & Partial & \checkmark & Partial & Partial & \checkmark \\

\addlinespace

\multicolumn{6}{l}{\textit{Interoperability}} \\
\midrule
Framework-independent representation       & -- & -- & -- & -- & \checkmark \\
Cross-framework portability                & -- & -- & Partial & -- & \checkmark \\

\bottomrule
\end{tabular}
\end{table*}

The comparison reveals that no existing domain individually provides all capabilities required to govern agentic systems throughout their lifecycle. SCM contributes mature configuration management principles but was not designed for dynamic AI artefacts. AI Governance establishes governance objectives without defining a configuration representation. LLMOps and AgentOps provide comprehensive operational observability but focus primarily on runtime execution. Agent frameworks implement execution semantics but remain tightly coupled to their internal abstractions.

These observations motivate the need for a unified reference model capable of bridging these complementary perspectives. Rather than replacing existing frameworks or operational platforms, such a model should provide a common representation of configuration artefacts, their relationships, their evolution, and their governance semantics. The remainder of this paper introduces Agentic Configuration Management (ACM) as such a framework-independent reference model.

\section{Design Principles}
\label{sec:design_principles}

The gap analysis presented in Section~\ref{sec:gap_analysis} highlights that existing approaches collectively address many aspects of agentic systems but fail to provide a unified representation capable of supporting configuration governance throughout the system lifecycle. Rather than extending an existing orchestration framework or operational platform, the objective of ACM is to define a framework-independent reference model that complements these ecosystems while remaining independent of their implementation choices.

To achieve this objective, the design of ACM is guided by the following principles.

\subsection{P1 -- Framework Independence}

The reference model shall remain independent of any specific orchestration framework, execution engine, or LLM provider. Agentic systems originating from heterogeneous ecosystems should be representable through a common abstraction without exposing framework-specific implementation details. This separation enables interoperability while preserving the semantics required for governance.

\subsection{P2 -- Explicit Configuration Representation}

Every artefact that influences the behaviour of an agentic system shall be represented explicitly as a configuration object. This includes not only prompts and models, but also agents, tools, workflows, policies, memories, runtime factories, evaluation artefacts, and other governance-related entities. Configuration should therefore be treated as a first-class concern rather than an implicit property of source code.

\subsection{P3 -- Immutable Configuration Evolution}

Configuration evolution shall be represented through immutable revisions rather than in-place modifications. Every released configuration must remain reproducible, auditable, and uniquely identifiable throughout its lifecycle. Historical configurations must therefore remain accessible independently of subsequent changes.

\subsection{P4 -- Separation of Configuration and Execution}

Configuration and execution represent complementary but fundamentally different concerns. The configuration model defines what the system is intended to execute, whereas runtime observations describe what actually occurred during execution. The reference model shall preserve this distinction while maintaining explicit traceability between configuration artefacts and runtime evidence.

\subsection{P5 -- Dependency-Aware Governance}

Configuration artefacts rarely evolve in isolation. Dependencies between artefacts influence system consistency, quality, assurance, and operational readiness. The reference model shall therefore explicitly represent dependency relationships and define governance semantics that preserve consistency across interconnected configuration objects.

\subsection{P6 -- Lifecycle-Oriented Governance}

Configuration management extends beyond version control. Each configuration artefact progresses through successive lifecycle states reflecting its maturity, validation status, and operational availability. The reference model shall support explicit lifecycle management together with deterministic governance rules governing state transitions and their consequences.

\subsection{P7 -- Explainable Governance Semantics}

Governance decisions should be transparent and reproducible. Rather than relying on implicit heuristics, the reference model shall define explicit semantics describing how assurance, eligibility, quality, impact, and other governance indicators are derived from the underlying configuration. Every computed governance state should therefore be explainable from observable configuration relationships and available evidence.

Together, these principles define the conceptual foundations of ACM. The following section introduces the reference model itself by describing its core concepts, structural relationships, and configuration abstractions.

\section{ACM Reference Model}
\label{sec:acm_reference_model}

This section introduces Agentic Configuration Management (ACM), a framework-independent reference model for representing, governing, and auditing the configuration of agentic systems. Rather than defining another execution framework, ACM provides a conceptual representation of the configuration itself, independently of the orchestration technology used to execute the system.

The reference model extends established principles from Software Configuration Management to the specific characteristics of agentic AI. It treats every artefact that contributes to the behaviour of an agentic system as an explicit configuration object and captures the structural relationships that exist among these artefacts. Configuration therefore becomes an explicit, versioned, and governable asset rather than an implicit consequence of implementation details.

The model is intentionally orthogonal to execution. Agent frameworks remain responsible for orchestration, scheduling, communication, and runtime execution, while ACM focuses on describing the configuration that is executed, the evolution of this configuration over time, and the evidence required to govern its lifecycle.

The ACM reference model is organized around three complementary concepts:

\begin{itemize}
    \item \textbf{Configuration artefacts}, representing the individual elements composing an agentic system;
    \item \textbf{Relationships}, describing the structural dependencies among these artefacts;
    \item \textbf{Configuration graphs}, providing multiple complementary views of the governed system.
\end{itemize}

Together, these concepts establish a normalized representation capable of supporting configuration governance independently of any particular implementation framework.

\subsection{Core Concepts}
\label{sec:core_concepts}

The fundamental abstraction of ACM is the \emph{Agentic Configuration Item (ACI)}. An ACI represents any identifiable artefact whose definition contributes to the observable behaviour of an agentic system and whose evolution should therefore be governed throughout its lifecycle.

Unlike traditional Software Configuration Management, where configuration items primarily consist of source files, libraries, or executable components, ACM extends the notion of configuration to encompass the heterogeneous artefacts specific to agentic systems. Prompts, agents, tools, models, workflows, policies, memories, evaluation profiles, runtime factories, and dynamically promoted artefacts are all treated as first-class configuration items.

Each ACI is uniquely identifiable, independently versioned, and immutable once released. Every modification results in a new revision while preserving the complete history of previous revisions. This approach guarantees reproducibility, auditability, and deterministic reconstruction of historical system configurations.

Although ACIs differ in their operational semantics, they share a common governance model based on explicit identity, immutable revisions, structural relationships, and lifecycle management. This common abstraction enables heterogeneous artefacts originating from different frameworks to be represented uniformly without constraining their internal implementation.

To preserve conceptual simplicity while remaining expressive, ACM organizes ACIs into three orthogonal families according to the role they play within the system.

\begin{enumerate}
    \item \textbf{Structural ACIs} describe the static configuration of the system, including agents, prompts, tools, models, workflows, memories, and governance policies.
    \item \textbf{Governance ACIs} represent the information required to validate, assess, and govern the configuration, such as evaluation profiles, evidence, approval records, and assurance artefacts.
    \item \textbf{Runtime ACIs} represent configuration objects originating from system execution and promoted for subsequent governance, including dynamically generated agents, prompts, workflows, or other persistent artefacts.
\end{enumerate}

These three families are intentionally orthogonal. A structural configuration object may be associated with governance artefacts through explicit relationships while remaining independent from runtime instances derived from its execution. Similarly, runtime-generated artefacts become governed configuration objects only after explicit promotion into the configuration repository.

This separation allows ACM to distinguish configuration from execution while preserving complete traceability between both perspectives.

\subsection{Configuration Item Model}
\label{sec:aci_model}

The Agentic Configuration Item (ACI) constitutes the fundamental entity of the ACM reference model. Every governable artefact participating in the definition, execution, validation, or governance of an agentic system is represented as an ACI. By treating heterogeneous artefacts through a common abstraction, ACM establishes a uniform representation independent of any orchestration framework.

Each ACI is uniquely identified by a stable identifier that remains constant throughout its lifetime. Its successive modifications are represented as immutable revisions, each identified by a revision identifier and associated with a cryptographic content digest. Consequently, the identity of an artefact is decoupled from its historical evolution, allowing both long-term traceability and deterministic reconstruction of previous configurations.

All ACI revisions share a common metadata structure composed of:

\begin{itemize}
    \item a persistent ACI identifier;
    \item an immutable revision identifier;
    \item a cryptographic digest of the revision content;
    \item creation timestamp;
    \item author or originating process;
    \item declared type;
    \item lifecycle metadata;
    \item governance metadata.
\end{itemize}

This common metadata layer provides a homogeneous governance interface regardless of the underlying artefact semantics.

ACM deliberately separates the notion of an artefact from the information describing its governance. The operational content of an ACI remains domain-specific, whereas lifecycle states, quality indicators, assurance evidence, dependencies, provenance, and revision history are represented through standardized governance attributes. This separation enables the reference model to evolve independently from implementation technologies.

The reference model distinguishes three orthogonal families of configuration items.

\paragraph{Structural ACIs}

Structural ACIs define the intended architecture of the agentic system. Typical examples include agents, prompts, workflows, tools, language models, memories, policies, knowledge bases, or external resources. Together they describe the configuration that is expected to be executed.

\paragraph{Governance ACIs}

Governance ACIs capture the evidence required to assess the configuration rather than its operational behaviour. They include evaluation reports, assurance artefacts, approval records, compliance evidence, validation procedures, and other governance documents that justify the operational readiness of structural artefacts.

\paragraph{Runtime ACIs}

Runtime ACIs represent configuration objects originating from execution. Autonomous systems may dynamically generate new prompts, instantiate specialized agents, synthesize workflows, or construct intermediate reasoning artefacts. Such objects initially belong to the execution domain and become governed configuration items only after explicit promotion into the ACM repository. This distinction preserves the separation between runtime behaviour and governed configuration while maintaining complete traceability between both worlds.

The classification of ACIs is orthogonal to their implementation technology. Consequently, identical conceptual artefacts originating from different orchestration frameworks can be represented using the same ACM abstractions.

\subsection{Relationships}
\label{sec:relationships}

Configuration artefacts rarely evolve independently. The behaviour of an agentic system emerges from the interactions among prompts, agents, workflows, tools, memories, policies, and runtime-generated objects. Representing individual artefacts without explicitly modelling their relationships would therefore provide only a partial description of the governed system.

For this reason, ACM models the complete configuration as a typed graph whose vertices correspond to ACI revisions and whose edges represent immutable semantic relationships. Relationships are first-class entities of the reference model and participate in governance in the same manner as configuration items.

Each relationship is directed, explicitly typed, and versioned together with the revisions it connects. Relationships therefore describe dependencies between immutable revisions rather than mutable configuration objects. This design guarantees deterministic reconstruction of historical configurations.

ACM defines a small set of generic relationship categories that remain independent of implementation frameworks.

\paragraph{Containment relationships}

Containment relationships describe structural composition. They indicate that one configuration item logically contains another as part of its definition. Examples include workflows containing agents, agents containing prompts, or policies containing evaluation criteria.

\paragraph{Dependency relationships}

Dependency relationships express that one configuration item requires another for its correct execution or validation. Dependencies are fundamental for impact analysis because modifications affecting one artefact may propagate to dependent configuration items.

\paragraph{Usage relationships}

Usage relationships capture operational interactions without implying ownership. Typical examples include an agent invoking a tool, a workflow referencing a language model, or a prompt accessing an external memory. Unlike containment, usage relationships allow artefacts to be shared among multiple configurations.

\paragraph{Governance relationships}

Governance relationships associate structural artefacts with the evidence required to justify their operational readiness. Evaluation reports, assurance artefacts, approval decisions, or compliance records are therefore explicitly linked to the configuration objects they govern.

\paragraph{Derivation relationships}

Derivation relationships preserve traceability between runtime execution and governed configuration. They connect dynamically generated artefacts with the structural configuration from which they originate, allowing autonomous evolution to remain auditable throughout the system lifecycle.

Because relationships are explicitly represented within the configuration graph, ACM supports deterministic impact analysis, dependency traversal, provenance reconstruction, and governance propagation without relying on framework-specific execution semantics. The meaning of each relationship type remains stable across implementation technologies, thereby enabling interoperability between heterogeneous orchestration frameworks.

\begin{table*}[t]
\centering
\caption{Normative relationship types defined by the ACM reference model.}
\label{tab:acm_relationships}
\small
\renewcommand{\arraystretch}{1.15}
\begin{tabular}{p{2.8cm}p{2.3cm}p{2.3cm}p{2.0cm}p{2.2cm}p{4.2cm}}
\toprule
\textbf{Relationship} &
\textbf{Source} &
\textbf{Target} &
\textbf{Cardinality} &
\textbf{Propagation Scope} &
\textbf{Semantics} \\
\midrule

\texttt{contains}
&
ACI
&
ACI
&
1:N
&
Recursive
&
Structural composition. The target is an integral part of the source configuration item.
\\

\texttt{depends\_on}
&
ACI
&
ACI
&
N:M
&
Recursive
&
Functional dependency. Changes to the target may affect the validity or readiness of the source.
\\

\texttt{uses}
&
ACI
&
ACI
&
N:M
&
Local
&
Operational usage without ownership. The target may be shared by multiple configuration items.
\\

\texttt{governed\_by}
&
Structural ACI
&
Governance ACI
&
N:M
&
None
&
Associates governance evidence (evaluation, approval, assurance, compliance) with a governed artefact.
\\

\texttt{derived\_from}
&
Runtime ACI
&
Structural ACI
&
1:N
&
None
&
Preserves provenance between runtime-generated artefacts and their originating configuration.
\\

\bottomrule
\end{tabular}
\end{table*}

\subsection{Configuration Graphs}
\label{sec:configuration_graphs}

Agentic systems cannot be adequately represented as a collection of independent configuration items. Their behaviour emerges from a network of structural dependencies, governance evidence, historical evolution, and runtime interactions. ACM therefore models the governed system through a set of complementary graphs, each describing a distinct aspect of the system while sharing the same underlying configuration items.

Rather than maintaining independent repositories, the ACM reference model defines multiple logical projections over a common configuration space. Each projection addresses a specific governance concern and collectively they provide a complete representation of the managed system.

\begin{figure*}[t]
	\centering
	\includegraphics[width=\textwidth]{figures/ACM_reference_model.png}
	\caption{The reference model of ACM.}
	\label{fig:ACM_reference}
\end{figure*}

\paragraph{Configuration Graph}

The Configuration Graph represents the intended structure of the agentic system. Vertices correspond to immutable ACI revisions, while edges represent the semantic relationships introduced in Section~\ref{sec:relationships}. This graph describes the configuration that is intended to be executed, independently of any runtime behaviour.

The Configuration Graph constitutes the primary representation manipulated by configuration management activities. It supports dependency analysis, impact assessment, baseline construction, and structural validation.

\paragraph{Evolution Graph}

The Evolution Graph captures the temporal evolution of every configuration item. Each revision is represented as an immutable node connected to its predecessor through explicit revision relationships. Unlike conventional version control systems, the Evolution Graph is defined at the configuration-model level rather than the source-code level, allowing heterogeneous artefacts to evolve under a common governance model.

Because revisions are immutable, every historical configuration remains reproducible. The complete evolution history of each ACI can therefore be reconstructed without ambiguity.

\paragraph{Runtime Graph}

The Runtime Graph represents the execution perspective of the system. It records runtime objects, execution traces, dynamically instantiated artefacts, and promoted runtime configuration items. Unlike the Configuration Graph, which represents the intended system, the Runtime Graph represents what actually occurred during execution.

Explicit derivation relationships preserve traceability between runtime artefacts and their originating structural configuration. Consequently, autonomous modifications performed during execution remain observable without modifying the original configuration.

\paragraph{Assurance Graph}

The Assurance Graph represents the governance evidence associated with configuration artefacts. Nodes correspond to governance ACIs such as evaluation reports, approval records, compliance evidence, validation results, or assurance artefacts, while edges associate this evidence with the structural configuration items they justify.

This representation separates governance evidence from operational configuration while preserving explicit traceability between both. Multiple independent evidence artefacts may therefore justify the operational readiness of a single configuration item.

Together, these four complementary graphs provide orthogonal views of the same governed system. Structural configuration, historical evolution, runtime execution, and governance evidence remain explicitly separated while sharing common identifiers and immutable revisions. This separation enables each concern to evolve independently without sacrificing overall consistency or traceability.

\subsection{Revision and Baseline Model}
\label{sec:revision_baseline}

Configuration governance requires not only explicit representation of configuration artefacts but also deterministic identification of their successive states. ACM therefore adopts an immutable revision model inspired by Software Configuration Management while extending it to heterogeneous agentic artefacts.

Every modification applied to a configuration item creates a new immutable revision. Existing revisions are never modified in place, ensuring that historical configurations remain permanently reproducible. Each revision is uniquely identified by its revision identifier and cryptographic content digest, allowing both semantic versioning and content verification.

The distinction between the persistent identity of an ACI and the identity of one of its revisions is fundamental. The persistent identifier denotes the conceptual configuration item, whereas revision identifiers describe successive immutable states of this item throughout its lifecycle. Relationships are established between revisions rather than mutable configuration objects, ensuring that historical dependency graphs remain stable over time.

A baseline represents a coherent and immutable snapshot of a configuration graph at a given point in time. Rather than storing duplicated artefacts, a baseline references the exact revisions that compose the governed configuration. Consequently, multiple baselines may share common revisions while remaining independently identifiable.

Because baselines are immutable, they provide deterministic reconstruction of previously validated configurations. This property supports reproducibility, auditing, certification, rollback, and longitudinal comparison of system evolution without ambiguity.

The reference model intentionally distinguishes configuration revisions from runtime observations. Runtime executions generate evidence associated with specific revisions but do not alter those revisions themselves. When execution results in the creation of persistent configuration artefacts, these artefacts are introduced as new ACIs through explicit derivation relationships and subsequently evolve under the same immutable revision model.

This approach preserves a clear separation between configuration evolution and execution history. The former describes intentional modifications to the governed system, whereas the latter documents operational observations generated during execution. Both perspectives remain explicitly connected through the graph model while preserving their respective semantics.

The immutable revision model established in this section provides the structural foundation required by the governance mechanisms introduced in the following section. Lifecycle transitions, assurance propagation, impact analysis, and configuration eligibility are all defined over immutable revisions and coherent baselines rather than mutable runtime objects.

\section{Governance Semantics}
\subsection{Lifecycle Model}
\label{sec:lifecycle}

The ACM reference model associates every immutable configuration revision with an explicit governance lifecycle. The lifecycle represents the governance status of a revision independently of its quality assessment, assurance evidence, or runtime execution. Consequently, lifecycle evolution captures governance decisions rather than structural modifications.

Let

\[
R
\]

denote the set of all immutable ACI revisions.

The lifecycle is defined as the function

\[
L:R\rightarrow S_L
\]

where

\[
S_L=
\{
Draft,
UnderReview,
Approved,
Operational,
Deprecated,
Archived
\}.
\]

For every revision

\[
r\in R,
\]

the value

\[
L(r)
\]

denotes its current lifecycle state.

The semantics of the lifecycle states are defined as follows.

\begin{itemize}

\item \textbf{Draft}

The revision has been created but has not yet entered the governance process.

\item \textbf{UnderReview}

The revision is being evaluated through validation, assessment, or organizational review.

\item \textbf{Approved}

The revision satisfies the required governance criteria and may participate in an operational baseline.

\item \textbf{Operational}

The revision belongs to an approved baseline and is eligible for operational deployment.

\item \textbf{Deprecated}

The revision remains valid for historical reproducibility but should no longer be selected for new deployments.

\item \textbf{Archived}

The revision is permanently preserved for auditability and historical reconstruction.

\end{itemize}

Because revisions are immutable, lifecycle changes never alter the configuration itself.

Let

\[
Content:R\rightarrow C
\]

represent the immutable content associated with every revision.

For every revision

\[
r\in R,
\]

the following invariant always holds:

\[
Content(r)=constant.
\]

Only the governance state

\[
L(r)
\]

may evolve during the lifetime of a revision.

This separation guarantees that governance decisions remain completely independent from configuration evolution. Any modification affecting the operational content of an artefact necessarily creates a new immutable revision

\[
r'
\neq
r,
\]

which starts a new lifecycle with

\[
L(r')=Draft.
\]

The lifecycle therefore governs immutable revisions rather than mutable configuration objects, ensuring deterministic governance and complete historical reproducibility.

\begin{figure*}[t]
	\centering
	\includegraphics[width=\textwidth]{figures/lifecyle_propagation_model_v2.png}
	\caption{Lifecyle and Sate Propagation Model.}
	\label{fig:ACM_propagation}
\end{figure*}


\subsection{Lifecycle State Machine}
\label{sec:lifecycle_state_machine}

The lifecycle model specifies the admissible governance states associated with immutable revisions. ACM further defines a deterministic state machine that constrains the possible transitions between these states.

Let

\[
T_L
\subseteq
S_L\times S_L
\]

denote the set of authorized lifecycle transitions.

The normative transition relation is defined as

\[
\begin{aligned}
T_L=\{
&
(Draft,UnderReview),\\
&
(UnderReview,Draft),\\
&
(UnderReview,Approved),\\
&
(Approved,Operational),\\
&
(Operational,Deprecated),\\
&
(Deprecated,Archived)
\}.
\end{aligned}
\]

A lifecycle transition is valid only if it belongs to this relation.

Formally,

\[
Transition(s_i,s_j)
=
\begin{cases}
true,&(s_i,s_j)\in T_L,\\
false,&\text{otherwise.}
\end{cases}
\]

For a revision

\[
r\in R,
\]

the application of a transition toward a new state

\[
s'
\]

is therefore defined as

\[
ApplyTransition(r,s')
=
\begin{cases}
L(r)\leftarrow s',
&
(L(r),s')\in T_L,\\
error,
&
\text{otherwise.}
\end{cases}
\]

Consequently, governance history is monotonic. Previously approved, operational, deprecated, or archived revisions cannot return to earlier governance states unless explicitly authorized by the transition relation.

Because ACM governs immutable revisions, prohibited transitions cannot be compensated by modifying an existing revision. Instead, any structural modification produces a new revision

\[
r'
\]

whose lifecycle is initialized as

\[
L(r')=Draft.
\]

This property guarantees that governance decisions remain permanently attached to the revision on which they were taken.

Although implementations may introduce additional intermediate states to support organization-specific governance processes, every compliant implementation shall preserve the transition semantics defined by

\[
T_L.
\]

The lifecycle state machine therefore provides a framework-independent governance semantics while allowing implementation-specific workflow refinements.

The state machine defined in this section constitutes the first formal governance mechanism of ACM. Subsequent sections define impact propagation, eligibility computation, and assurance evaluation over this lifecycle foundation.

\subsection{Impact Semantics}
\label{sec:impact}

Configuration changes rarely affect a single artefact in isolation. Because configuration items are interconnected through explicit semantic relationships, modifications may influence dependent artefacts, governance evidence, operational readiness, or deployed baselines. ACM therefore defines impact as a graph-based property evaluated over immutable configuration revisions.

Let

\[
G=(V,E)
\]

denote the Configuration Graph, where

\[
V=\{r_1,r_2,\ldots,r_n\}
\]

is the set of immutable ACI revisions and

\[
E\subseteq V\times T\times V
\]

is the set of typed relationships introduced in Section~\ref{sec:relationships}, with

\[
T=
\{
contains,
depends\_on,
uses,
governed\_by,
derived\_from
\}.
\]

Only relationships participating in governance propagation are considered during impact computation. Let

\[
T_p=
\{
contains,
depends\_on,
uses
\}
\]

denote this propagation subset.

Given a modified revision \(r\), the direct impact set is defined as

\[
I_1(r)
=
\{
v\in V
\mid
(r,t,v)\in E
\land
t\in T_p
\}.
\]

The transitive impact set is recursively computed as

\[
I^{+}(r)
=
I_1(r)
\cup
\bigcup_{v\in I_1(r)} I^{+}(v).
\]

The complete impact closure therefore corresponds to the transitive closure of propagation relationships.

Governance evidence does not itself participate in impact propagation. Instead, governance artefacts remain attached to impacted revisions through \texttt{governed\_by} relationships and are subsequently re-evaluated if required.

Because ACM operates on immutable revisions, impact analysis is deterministic. The same configuration graph always produces the same impacted set independently of execution order or implementation framework.

The impact semantics constitute the foundation upon which eligibility and assurance propagation are subsequently evaluated.

\subsection{Eligibility Semantics}
\label{sec:eligibility}

Lifecycle progression alone does not guarantee that a configuration revision may be deployed. Operational readiness additionally depends on governance constraints such as approval status, assurance evidence, dependency validity, and organizational policies. ACM captures these conditions through the notion of eligibility.

Eligibility is evaluated independently for every immutable revision.

Let

\[
Eligible(r)\in\{true,false\}
\]

denote the eligibility predicate associated with revision \(r\).

A revision is eligible only if all required governance conditions are simultaneously satisfied.

Formally,

\[
Eligible(r)
=
L(r)
\land
A(r)
\land
D(r)
\land
P(r),
\]

where

\begin{itemize}

\item \(L(r)\) denotes lifecycle validity,

\item \(A(r)\) denotes assurance validity,

\item \(D(r)\) denotes dependency validity,

\item \(P(r)\) denotes policy compliance.

\end{itemize}

Lifecycle validity is defined as

\[
L(r)
=
\big(
State(r)=Operational
\big).
\]

Dependency validity requires that every propagated dependency remains eligible.

Let

\[
Dep(r)
=
\{
v
\mid
(r,depends\_on,v)\in E
\}.
\]

Then

\[
D(r)
=
\bigwedge_{v\in Dep(r)}
Eligible(v).
\]

Consequently, eligibility naturally propagates along dependency relationships. The loss of eligibility of one configuration item may therefore invalidate all dependent revisions.

Policy compliance remains implementation-independent. ACM defines only the evaluation interface, while concrete policy rules are supplied by compliant implementations.

Eligibility is always computed rather than manually assigned. It therefore represents a derived governance property whose value is entirely determined by the current state of the configuration graph.

\subsection{Assurance Semantics}
\label{sec:assurance}

Configuration governance requires explicit evidence demonstrating that configuration artefacts satisfy their expected quality and operational requirements. Rather than associating a single approval status with each revision, ACM models assurance as an explicit graph of evidence connected to immutable configuration items.

Each structural revision declares the assurance dimensions required for operational use.

Let

\[
R(r)
=
\{
d_1,d_2,\ldots,d_n
\}
\]

denote the required assurance dimensions for revision \(r\).

Typical dimensions include functional validation, safety assessment, security evaluation, regulatory compliance, performance testing, or organizational approval.

Evidence is represented by Governance ACIs connected through \texttt{governed\_by} relationships.

For each assurance dimension \(d\), let

\[
Evidence(r,d)
=
\{
e_1,\ldots,e_k
\}
\]

denote the available evidence associated with revision \(r\).

An assurance dimension is considered satisfied whenever at least one valid evidence artefact covers the corresponding requirement.

Formally,

\[
Satisfied(r,d)
=
\exists e
\in
Evidence(r,d)
:
Valid(e).
\]

Overall assurance validity becomes

\[
A(r)
=
\bigwedge_{d\in R(r)}
Satisfied(r,d).
\]

A revision therefore satisfies its assurance requirements only if every declared assurance dimension is covered by valid governance evidence.

Assurance propagation follows dependency relationships.

For every dependency

\[
(r,depends\_on,v),
\]

the effective assurance of \(r\) is bounded by the weakest assurance of its dependencies.

This property is expressed through

\[
EffectiveAssurance(r)
=
\min
\left(
A(r),
\min_{v\in Dep(r)}
EffectiveAssurance(v)
\right).
\]

Consequently, a configuration cannot exhibit a higher effective assurance level than any of the configuration items upon which it depends.

Separating assurance evidence from structural configuration provides several advantages. Evidence may evolve independently from configuration revisions, multiple independent evaluations may coexist, and assurance remains fully traceable throughout the evolution of the configuration graph.

The assurance semantics complete the governance foundations of ACM. Together with lifecycle management, impact analysis, and eligibility computation, they establish a deterministic governance model that remains independent of any particular agent orchestration framework.


\subsection{Propagation Semantics}
\label{sec:propagation}

The governance attributes introduced in the previous sections are not evaluated independently. Because configuration revisions are interconnected through explicit semantic relationships, governance properties propagate over the configuration graph according to deterministic rules.

Let

\[
G=(V,E)
\]

be the Configuration Graph introduced in Section~\ref{sec:configuration_graphs}, where

\[
V=R.
\]

Only propagation relationships participate in governance propagation:

\[
T_p=
\{
contains,
depends\_on,
uses
\}.
\]

The neighbourhood of a revision is therefore defined as

\[
N(r)
=
\{
v
\in
R
\mid
(r,t,v)\in E
\land
t\in T_p
\}.
\]

Governance propagation is defined as the least fixed point of repeated local evaluations.

Let

\[
X(r)
\]

represent any propagated governance attribute (eligibility, assurance, quality, or future governance metrics).

The propagation operator is defined by

\[
X_{k+1}(r)
=
f
\left(
X_k(r),
\{
X_k(v)
\mid
v\in N(r)
\}
\right),
\]

where

\[
f
\]

is the governance-specific aggregation function.

The propagation process converges when

\[
X_{k+1}(r)
=
X_k(r)
\qquad
\forall r\in R.
\]

Because the configuration graph is finite and governance attributes evolve monotonically, the propagation process always converges toward a unique fixed point.

Consequently, every governance computation performed by ACM is deterministic and independent of traversal order.

The concrete aggregation function depends on the propagated attribute. For example, eligibility uses logical conjunction, assurance uses weakest-dependency semantics, while future governance dimensions may define alternative aggregation operators without modifying the propagation mechanism itself.

\subsection{Runtime Semantics}
\label{sec:runtime_semantics}

The Configuration Graph represents the intended configuration of an agentic system. Runtime execution, however, produces additional information that does not necessarily modify the governed configuration. ACM therefore explicitly separates configuration semantics from execution semantics.

Let

\[
E_x
\]

denote the set of runtime execution events.

Each event

\[
e
\in
E_x
\]

is associated with exactly one immutable revision

\[
Origin(e)
\in
R.
\]

Runtime observations are therefore represented by the mapping

\[
Origin:E_x\rightarrow R.
\]

Execution events never modify configuration revisions.

Formally,

\[
\forall
e
\in
E_x,
\qquad
Content(Origin(e))
=
constant.
\]

Instead, runtime observations extend the Runtime Graph while preserving complete traceability toward the originating configuration.

Execution semantics therefore operate over

\[
G_R=(E_x,E_R),
\]

where

\[
E_R
\]

represents runtime relationships such as execution order, event causality, or promoted runtime artefacts.

The Runtime Graph and Configuration Graph remain explicitly connected through the mapping

\[
Origin.
\]

This separation guarantees that execution history never alters the governed configuration while preserving complete provenance between design-time and runtime representations.

\subsection{Dynamic Agent Semantics}
\label{sec:dynamic_agents}

Autonomous agentic systems may dynamically generate new prompts, instantiate specialized agents, synthesize workflows, or create intermediate reasoning artefacts during execution.

ACM distinguishes such runtime artefacts from governed configuration items.

Let

\[
D
\subseteq
E_x
\]

denote the set of dynamically generated runtime artefacts.

Initially,

\[
D
\cap
R
=
\varnothing.
\]

Dynamic artefacts therefore belong exclusively to the Runtime Graph.

A runtime artefact becomes governed only through an explicit promotion operation.

Let

\[
Promote
:
D
\rightarrow
R
\]

be the promotion function.

For every promoted artefact

\[
d
\in
D,
\]

there exists a newly created immutable revision

\[
r
=
Promote(d).
\]

This revision satisfies

\[
L(r)=Draft.
\]

Furthermore,

\[
(r,derived\_from,d)
\]

is added to the Configuration Graph, preserving complete provenance between runtime generation and governed configuration.

Promotion never modifies the original runtime artefact.

Instead,

\[
Promote(d)
\neq
d.
\]

The promoted revision subsequently evolves according to the standard lifecycle, eligibility, assurance, and propagation semantics defined by ACM.

Consequently, autonomous runtime evolution remains fully compatible with deterministic configuration governance without violating the immutability principles established by the reference model.

\subsection{Governance Semantics Summary}
\label{sec:governance_summary}

This section established the formal governance semantics of the ACM reference model.

Starting from the lifecycle function

\[
L:R\rightarrow S_L,
\]

ACM defines deterministic state transitions through the lifecycle state machine, impact propagation over the configuration graph, eligibility computation, assurance evaluation, governance propagation, runtime semantics, and dynamic agent promotion.

Together, these semantics transform the structural model introduced in Section~4 into a complete governance model. Every governance property is computed over immutable revisions connected through explicit semantic relationships, ensuring deterministic behaviour, complete traceability, and reproducibility across heterogeneous agent orchestration frameworks.

The resulting governance model is intentionally independent of implementation technologies. Framework-specific orchestration mechanisms therefore become concrete realizations of a common mathematical semantics rather than independent governance models.

The following section demonstrates that these semantics can be operationalized through a reference implementation and evaluated across representative agentic frameworks while preserving framework independence.

\section{Operationalization}
\label{sec:operationalization}
\subsection{Reference Architecture}
\label{sec:reference_architecture}

The ACM reference model is intentionally independent of any execution framework. Rather than embedding orchestration capabilities, the reference implementation operationalizes the conceptual model through a layered architecture that separates framework-specific extraction from framework-independent governance.

Figure~\ref{fig:projection_pipeline} illustrates the overall architecture. Native agentic systems are first interpreted by framework adapters that expose their configuration through a common representation. Once normalized, the resulting ACM objects are processed exclusively by the framework-independent governance kernel.

The reference implementation is organized around four logical layers.

\paragraph{Framework adapters.}

Framework adapters constitute the only framework-dependent components of the implementation. Their responsibility is to extract the declarative configuration exposed by a supported framework and translate it into ACM entities without modifying its semantics. The current implementation provides adapters for LangGraph and CrewAI, allowing heterogeneous workflow descriptions to be projected onto a common configuration model. Adapter outputs are expressed exclusively in terms of ACM concepts and therefore isolate the governance kernel from framework-specific APIs. :contentReference[oaicite:2]{index=2}

\paragraph{Normalization layer.}

The normalization layer converts extracted artifacts into typed ACM objects. Every configuration element is represented as an Agentic Configuration Item (ACI) associated with immutable identifiers, explicit revisions, canonical content digests and normalized relationships. During this phase, framework-specific concepts are preserved whenever possible while being mapped onto the ACM metamodel. The objective is not to reinterpret the source framework but to construct a configuration graph that can subsequently be validated and governed using framework-independent semantics. :contentReference[oaicite:3]{index=3}

\paragraph{Governance kernel.}

The governance kernel implements the operational semantics defined in Sections~4 and~5. It validates structural consistency, evaluates configuration invariants, computes derived states, and applies deterministic propagation rules over the configuration graph. The kernel manipulates only normalized ACM objects and is therefore independent of the framework from which the configuration originated.

State computation follows the separation introduced by ACM between intrinsic states (lifecycle, quality and assurance) and computed states (impact, eligibility and effective states). These values are derived by a deterministic fixed-point computation whose result depends solely on the normalized configuration, the associated evidence and the applicable propagation policy. Consequently, identical ACM inputs produce identical governance results regardless of their original framework representation. :contentReference[oaicite:4]{index=4} :contentReference[oaicite:5]{index=5}

\paragraph{Runtime interface.}

Execution remains external to ACM. Runtime frameworks interact with the governance kernel through a dedicated runtime interface that produces normalized runtime signals instead of exposing native execution objects directly. This design preserves the separation between configuration governance and execution while enabling replay, provenance reconstruction and drift analysis from normalized runtime events. The governance kernel therefore reasons over framework-independent runtime representations rather than execution-specific APIs. :contentReference[oaicite:6]{index=6}

The resulting architecture deliberately minimizes coupling between execution frameworks and governance mechanisms. Framework-specific evolution is confined to adapters, whereas configuration validation, lifecycle management, provenance, assurance computation and state propagation remain entirely implemented within the ACM kernel. This separation allows new execution frameworks to be integrated by implementing an additional projection layer without modifying the governance semantics defined by ACM.

\subsection{Projection Pipeline}
\label{sec:projection_pipeline}

The ACM reference implementation operationalizes framework-independent governance through a deterministic projection pipeline that transforms native framework configurations into normalized ACM representations. Rather than executing or modifying workflows, the pipeline extracts their declared configuration, preserves their structural semantics whenever possible, and produces a canonical representation suitable for validation and governance.

As illustrated in Figure~\ref{fig:projection_pipeline}, the projection process is intentionally decomposed into independent stages. Each stage has a single responsibility, allowing framework-specific concerns to remain isolated from the governance kernel while ensuring that every transformation remains reproducible and auditable.

\subsection{Framework Adapters}
\label{sec:framework_adapters}

The ACM reference model is intentionally independent of any particular orchestration framework. Rather than defining native execution semantics, it specifies a normalized configuration model together with a deterministic governance semantics. Framework adapters bridge the gap between native framework representations and this common model by extracting configuration artifacts, preserving their structural relationships whenever possible, and producing equivalent ACM representations.

Each adapter follows the projection pipeline presented in the previous section. Native configuration objects are first identified and resolved, then translated into Agentic Configuration Items (ACIs) and their associated relationships. The resulting ACM graph is subsequently validated and governed using the framework-independent algorithms described in Section~\ref{sec:governance_algorithms}. Consequently, all governance operations are performed on the normalized ACM representation rather than on framework-specific data structures.

The reference implementation currently provides adapters for LangGraph and CrewAI. These two frameworks were selected because they represent two complementary orchestration paradigms: explicit stateful execution graphs and role-oriented collaborative orchestration. Although their internal abstractions differ substantially, both can be projected into the same ACM representation, providing initial evidence that the reference model is portable across heterogeneous agentic frameworks.

\subsubsection{LangGraph}

LangGraph represents agentic applications as explicit directed state graphs in which nodes encapsulate executable functions or agents and edges define deterministic routing conditions. Shared state is propagated throughout the graph and updated after each node execution, making the execution topology directly observable.

The LangGraph adapter exploits this explicit representation to reconstruct the declared configuration with minimal interpretation. Graph nodes are projected as workflow nodes referencing \texttt{AgentDefinition} instances, while execution edges are translated into workflow transitions preserving their declared control-flow semantics. Agent definitions, prompts, model configurations and tool bindings are extracted from the corresponding LangGraph objects and normalized as ACM configuration items. Runtime checkpoints and execution traces are converted into normalized runtime events without modifying the baseline configuration.

Because LangGraph explicitly exposes both topology and execution state, the projection preserves most structural relationships required by the ACM reference model. Information loss is therefore limited primarily to implementation-specific execution metadata that has no equivalent in the framework-independent representation.

\subsubsection{CrewAI}

CrewAI adopts a higher-level orchestration model centred on collaborating agents, tasks, crews and, more recently, flows. Unlike LangGraph, the execution graph is generally not represented explicitly as a static graph but emerges from the interactions between agents, task assignments and orchestration primitives.

The CrewAI adapter reconstructs an equivalent ACM configuration by identifying the participating agents, their associated tasks, prompt definitions, models, tools and orchestration metadata. Task dependencies and flow transitions are translated into normalized workflow relationships whenever an explicit correspondence exists. When execution semantics are implicit or delegated to the framework runtime, the adapter records the corresponding approximation while preserving complete traceability between the extracted ACM objects and their native CrewAI sources.

This preservation-based extraction strategy ensures that no unsupported semantics are artificially introduced into the ACM model. Instead, any information that cannot be represented exactly is explicitly identified as an approximation. This design preserves the determinism of the governance process while making the boundaries of the projection transparent to users and evaluators.

\subsection{Governance Algorithms}
\label{sec:governance_algorithms}

The ACM reference model is operationalized through four deterministic algorithms that transform framework-specific configurations into governed ACM representations. Together, these algorithms implement the complete governance workflow, from configuration extraction to runtime state reconstruction. Their execution order is illustrated in Figure~\ref{fig:projection_pipeline}.

The algorithms are intentionally independent of any particular orchestration framework. Framework adapters are responsible only for extracting native configuration artifacts. All subsequent validation, state computation and governance decisions are performed exclusively on the normalized ACM representation.

\paragraph{Algorithm 1. Configuration Projection}

The projection algorithm transforms a native framework configuration into an ACM Configuration Graph. It identifies the configuration artifacts exposed by the framework, creates the corresponding Agentic Configuration Items (ACIs), resolves their references and reconstructs the declared dependency graph.

The algorithm follows a preservation-based strategy: whenever an exact semantic correspondence exists, the native artifact is translated directly into its ACM equivalent. When no exact mapping exists, the adapter records the approximation explicitly rather than introducing inferred semantics. This guarantees that the resulting ACM graph faithfully represents the source configuration while preserving complete traceability to the original framework objects.

\begin{algorithm}[H]
\caption{Configuration Projection}
\label{alg:projection}
\KwIn{Native framework configuration $F$}
\KwOut{Configuration graph $G_C$}

Initialize empty configuration graph $G_C$\;

\ForEach{configuration artifact $a \in F$}{
    create corresponding ACI\;
    normalize metadata and identifiers\;
    compute canonical revision identifier and digest\;
    add ACI to $G_C$\;
}

\ForEach{native relationship $r$}{
    resolve referenced ACIs\;
    create normalized ACM relationship\;
}

validate referential consistency\;

\Return{$G_C$}
\end{algorithm}

\paragraph{Algorithm 2. Structural Validation}

Once projected, the configuration graph is validated against the ACM specification. Structural validation verifies schema conformance, identifier uniqueness, referential integrity, dependency consistency and policy constraints before any governance computation is performed.

Validation intentionally separates structural correctness from governance decisions. The algorithm reports every detected inconsistency but does not modify the configuration graph. This separation preserves reproducibility while ensuring that governance algorithms always operate on a well-defined representation.

\begin{algorithm}[H]
\caption{Structural Validation}
\label{alg:validation}
\KwIn{Configuration graph $G_C$}
\KwOut{Validation report $R$}

initialize empty report $R$\;

validate schema constraints\;
validate unique identifiers\;
resolve all references\;
validate relationship cardinalities\;
validate lifecycle constraints\;
validate applicable policies\;

record every detected violation in $R$\;

\Return{$R$}
\end{algorithm}

\paragraph{Algorithm 3. Fixed-Point Governance Propagation}

Governance decisions are computed by iterative propagation over the configuration graph. Starting from the intrinsic states declared by individual ACIs, the algorithm repeatedly computes impact, eligibility, effective quality and effective assurance until no further state changes occur.

Unlike hierarchical propagation, the algorithm never overwrites intrinsic states. Instead, it derives computed states while preserving the original lifecycle, quality and assurance assigned to each revision. This distinction ensures that governance remains fully explainable and reproducible.

Because each iteration is monotonic with respect to the propagation rules, the computation converges toward a stable fixed point.

\begin{algorithm}[H]
\caption{Fixed-Point Governance Propagation}
\label{alg:propagation}
\KwIn{Configuration graph $G_C$}
\KwOut{Governed configuration graph}

initialize computed states from intrinsic states\;

\Repeat{no state changes}{
    \ForEach{ACI}{
        compute dependency impact\;
        compute effective assurance\;
        compute effective quality\;
        compute eligibility state\;
    }
}

\Return{governed configuration graph}
\end{algorithm}

\paragraph{Algorithm 4. Runtime Integration}

Runtime integration combines an immutable baseline with normalized runtime events to reconstruct the observed execution state. Instead of modifying the released configuration directly, runtime mutations are recorded as immutable events that can be replayed deterministically.

This event-sourcing approach preserves the distinction between declared configuration and observed execution while enabling configuration drift detection and complete execution reconstruction.

\begin{algorithm}[H]
\caption{Runtime Integration}
\label{alg:runtime}
\KwIn{Release baseline $B$, runtime event stream $E$}
\KwOut{Runtime graph $G_R$}

initialize runtime graph from baseline $B$\;

\ForEach{event $e \in E$}{
    validate event\;
    apply runtime mutation\;
    update runtime state\;
}

compute runtime governance status\;

\Return{$G_R$}
\end{algorithm}

\paragraph{Complexity Analysis.}

Let $N$ denote the number of ACIs, $M$ the number of relationships, $K$ the number of propagation iterations required to reach a fixed point, and $R$ the number of runtime events.

The projection algorithm performs a single traversal of the configuration artifacts and their relationships, yielding a time complexity of $O(N+M)$. Structural validation also executes in $O(N+M)$, assuming constant-time identifier lookups through indexed data structures.

Governance propagation dominates the computational cost. Each iteration evaluates every configuration item and its incoming dependencies, leading to a worst-case complexity of $O(K(N+M))$. In practice, convergence is typically achieved after a small number of iterations because state transitions are monotonic and the propagation rules are deterministic.

Runtime integration processes the event log sequentially, resulting in a complexity of $O(R)$. Since the baseline remains immutable, replay complexity depends only on the number of recorded runtime events and is independent of the original framework representation.

The reference implementation further reduces practical execution time by maintaining indexed relationships and deterministic traversal orders, avoiding repeated graph searches during successive propagation iterations.

\subsection{Reference Implementation}
\label{sec:reference_implementation}

To demonstrate the practical applicability of the proposed reference model, ACM has been implemented as a lightweight Python reference implementation. The objective of this prototype is not to provide a production-ready governance platform, but to demonstrate that the concepts introduced in Sections~\ref{sec:reference_model} and~\ref{sec:governance_semantics} can be implemented using deterministic algorithms and conventional software engineering techniques.

The implementation follows the architecture introduced in Section~\ref{sec:reference_architecture}. Framework adapters first extract native configurations and project them into the ACM metamodel. The governance kernel subsequently performs structural validation, state computation and runtime integration independently of the originating framework. This strict separation ensures that governance semantics remain entirely framework-independent.

The implementation relies on Pydantic models to define the normative structure of every Agentic Configuration Item (ACI). Schema validation is therefore performed automatically during object construction, guaranteeing structural consistency before governance algorithms are executed. Each ACI is represented as a strongly typed immutable revision identified by a logical identifier, a unique revision identifier and a canonical cryptographic digest.

Canonical serialization plays a central role in ensuring reproducibility. Every configuration object is serialized using a deterministic representation prior to digest computation, guaranteeing that semantically equivalent configurations always produce identical content identifiers regardless of platform or execution environment. Released baselines are therefore content-addressable and tamper-evident, allowing reproducible configuration snapshots to be reconstructed from their recorded revisions.

Structural validation is implemented as an independent validation engine executed immediately after projection. The engine verifies schema conformance, referential integrity, relationship consistency, lifecycle constraints and applicable governance policies before any propagation takes place. Validation reports remain separate from governance decisions in order to preserve a clear distinction between structural correctness and computed operational states.

Governance decisions are produced by a deterministic propagation engine implementing the fixed-point algorithm described previously. Starting from the intrinsic states associated with individual revisions, the engine iteratively computes impact, eligibility, effective quality and effective assurance until convergence is reached. Intrinsic states are never modified during this process; only derived governance states are updated, preserving complete traceability between declared configuration and computed governance outcomes.

Runtime observations are managed through an append-only event log following an event-sourcing approach. Instead of modifying released baselines, runtime events record the mutations observed during execution. Replaying this normalized event sequence reconstructs the Runtime Graph deterministically while preserving the original immutable configuration. This separation enables configuration drift analysis, execution replay and post-execution auditing without altering the published baseline.

Framework interoperability is achieved through dedicated adapters that isolate all framework-specific extraction logic. The current implementation includes adapters for LangGraph and CrewAI, demonstrating that distinct orchestration paradigms can be normalized into the same ACM representation. Because governance algorithms operate exclusively on normalized ACIs, extending ACM to additional frameworks requires only the implementation of a new projection adapter, while the governance kernel itself remains unchanged.

The reference implementation was intentionally designed as a modular architecture composed of loosely coupled components. Validation rules, propagation policies, runtime adapters and future governance mechanisms can therefore evolve independently without modifying the core metamodel. This extensibility supports the long-term objective of ACM as a framework-independent reference model rather than a framework-specific execution platform.

\section{Evaluation}
\label{sec:evaluation}
\subsection{Evaluation Objectives}
\label{sec:evaluation_objectives}

The objective of this evaluation is not to benchmark an orchestration framework or to compare execution performance. Instead, it assesses whether the proposed ACM reference model satisfies the design goals introduced in Section~\ref{sec:goals} and whether its reference implementation faithfully operationalizes the governance semantics defined in Sections~\ref{sec:model} and~\ref{sec:governance_semantics}.

Because ACM is proposed as a framework-independent reference model, the evaluation focuses on four complementary dimensions: (i) the expressiveness of the model for representing governed agentic configurations, (ii) the consistency of its governance semantics and executable invariants, (iii) the portability of the projection process across heterogeneous orchestration frameworks, and (iv) the feasibility of implementing the model using deterministic software engineering techniques.

These objectives are formalized through the following research questions.

\begin{description}[leftmargin=1.2cm,style=nextline]

\item[\textbf{RQ1. Model Expressiveness}]
Can the ACM reference model represent and govern the configuration concepts required for agentic systems, including versioned configuration items, lifecycle states, baselines, assurance evidence, runtime events, and their relationships?

\item[\textbf{RQ2. Governance Consistency}]
Do the validation, propagation and replay mechanisms consistently enforce the normative semantics defined by the ACM specification through executable invariants and deterministic state computation?

\item[\textbf{RQ3. Cross-Framework Portability}]
Can heterogeneous agentic workflows be projected into a common ACM representation while preserving their configuration semantics and explicitly characterizing any approximation or unsupported information?

\item[\textbf{RQ4. Implementation Feasibility}]
Can the ACM reference model be implemented as a lightweight and deterministic software layer using conventional engineering techniques without depending on a specific orchestration framework?

\end{description}

The remainder of this section addresses these research questions through a structured evaluation methodology. Section~\ref{sec:experimental_protocol} first describes the experimental protocol and evaluation assets. The normative scenarios are then used to validate the governance semantics of the model, followed by a cross-framework evaluation of LangGraph and CrewAI projections. The analysis subsequently examines information preservation during projection before assessing the reference implementation and discussing the implications and limitations of the obtained results.

\subsection{Experimental Protocol}
\label{sec:experimental_protocol}

The evaluation protocol was designed to assess the ACM reference model rather than the software implementation itself. Consequently, the experiments focus on verifying that the proposed governance semantics can be operationalized consistently, that representative agentic configurations can be projected into the ACM representation, and that the resulting governed configurations satisfy the normative properties defined throughout this paper.

The experimental campaign relies on the Python reference implementation described in Section~\ref{sec:reference_implementation}. This implementation serves exclusively as an executable realization of the ACM specification and provides a deterministic environment for validating the proposed algorithms and governance rules. The evaluation does not compare ACM against alternative orchestration frameworks or governance solutions, nor does it attempt to benchmark execution performance.

The complete evaluation environment consists of the ACM reference implementation executed on a Linux \texttt{x86\_64} platform with Python~3.13.14 and the current LangGraph and CrewAI adapters enabled. The latest complete validation report executes 277 automated tests without failures or skipped tests. In addition, eleven propagation scenarios are represented as declarative YAML fixtures executed through a common validation harness in order to verify the determinism of the propagation engine independently of the Python test implementation. :contentReference[oaicite:2]{index=2}

The experimental protocol is organized around four complementary evaluation stages corresponding to the research questions introduced previously.

\begin{enumerate}[leftmargin=*]
    \item \textbf{Normative scenario validation.} The ACM governance semantics are evaluated using a catalogue of normative scenarios covering the principal invariants defined by the specification. Each scenario represents a representative configuration or runtime situation together with its expected governance outcome.

    \item \textbf{Cross-framework projection.} Representative workflows extracted from LangGraph and CrewAI are projected into the ACM metamodel using the framework adapters presented in Section~\ref{sec:framework_adapters}. The objective is to assess whether heterogeneous orchestration models can be normalized into a common configuration representation.

    \item \textbf{Information preservation analysis.} The normalized ACM representations are compared with manually established reference representations in order to classify extracted information as preserved, approximated, or unsupported within the normative scope of ACM. This analysis characterizes the fidelity of framework projection rather than attempting to measure semantic equivalence between orchestration frameworks.

    \item \textbf{Reference implementation assessment.} Finally, the consistency of the executable specification is assessed through the complete automated validation suite, including deterministic propagation tests, structural validation, runtime replay, and regression scenarios.
\end{enumerate}

The normative evaluation assets comprise twenty-seven scenarios grouped into five categories: (i) configuration and baseline integrity, (ii) propagation and assurance, (iii) runtime replay and configuration drift, (iv) dynamic agents and permission management, and (v) portability and robustness. Each scenario is associated with one or more executable tests and collectively provides coverage of the principal governance mechanisms introduced by the ACM specification. :contentReference[oaicite:3]{index=3}

To ensure reproducibility, all experiments are executed using deterministic configuration revisions, canonical serialization, immutable baselines, and normalized runtime events. Framework-specific extraction is confined to the projection phase; every subsequent validation, propagation, replay and governance computation operates exclusively on the normalized ACM representation. This separation ensures that the reported results evaluate the reference model itself rather than the internal behavior of the supported orchestration frameworks.

\subsection{Normative Scenario Validation}
\label{sec:normative_validation}

The primary objective of the evaluation is to determine whether the ACM reference implementation satisfies the normative semantics introduced in Sections~\ref{sec:reference_model} and~\ref{sec:governance_semantics}. Rather than validating individual software functions in isolation, the evaluation is organized around a catalogue of executable normative scenarios. Each scenario represents a representative configuration or runtime situation together with the expected governance behaviour prescribed by the ACM specification.

The complete evaluation suite contains twenty-seven scenarios grouped into five complementary categories. Collectively, these scenarios exercise the principal concepts introduced by ACM, including immutable revisions, baseline integrity, dependency propagation, assurance composition, runtime replay, configuration drift, dynamic agents, permission ceilings, portability across orchestration frameworks, determinism and local scalability. Every scenario is associated with one or more automated tests, providing an executable interpretation of the specification. The complete validation campaign reports 277 executed tests with no failures and no skipped tests. :contentReference[oaicite:0]{index=0}

\begin{table}[t]
\centering
\caption{Normative scenario categories.}
\label{tab:scenario_categories}
\begin{tabularx}{\linewidth}{lcl}
\toprule
Group & Scenarios & Evaluation scope \\
\midrule
A & S01--S04 &
Configuration and baseline integrity \\
B & S05--S11 &
Propagation, assurance and governance semantics \\
C & S12--S17 &
Runtime replay and configuration drift \\
D & S18--S21 &
Dynamic agents and permission management \\
E & S22--S27 &
Portability, robustness and local scalability \\
\bottomrule
\end{tabularx}
\end{table}

Table~\ref{tab:scenario_results} summarizes the outcome of the complete scenario campaign. Twenty-four scenarios satisfy both the specification and the implementation without qualification. Three runtime-related scenarios are classified as \texttt{pass\_with\_deviation}. In these cases, all executable tests succeed, but the implementation intentionally derives a more detailed conformity interpretation than originally prescribed by the evaluation plan. These deviations correspond to modelling choices documented by the implementation and do not represent validation failures. :contentReference[oaicite:1]{index=1}

\begin{table*}[t]
\centering
\caption{Summary of normative scenario validation.}
\label{tab:scenario_results}
\begin{tabularx}{\textwidth}{lccccX}
\toprule
Group &
Scenarios &
Pass &
Pass with deviation &
Fail &
Primary validation objective \\
\midrule
A &
4 &
4 &
0 &
0 &
Revision identity, baseline integrity, referential consistency and content digests \\

B &
7 &
7 &
0 &
0 &
Propagation semantics, assurance composition, quality propagation and eligibility \\

C &
6 &
3 &
3 &
0 &
Runtime replay, state transitions and configuration drift \\

D &
4 &
4 &
0 &
0 &
Dynamic agents, permission ceilings, behavioural overrides and promotion \\

E &
6 &
6 &
0 &
0 &
Framework portability, ordering invariance, dependency cycles and local scalability \\

\midrule
\textbf{Total} &
\textbf{27} &
\textbf{24} &
\textbf{3} &
\textbf{0} &
277 automated tests, no failures and no skipped tests \\
\bottomrule
\end{tabularx}
\end{table*}

The three documented deviations concern scenarios S14--S16, all belonging to the runtime and configuration-drift category. In these cases, the governance engine reports a compact internal drift state while the conformity layer derives a more expressive classification intended for reporting purposes. Consequently, concepts such as \emph{declared extension}, \emph{untraceable instance} and configuration mismatch are represented as derived conformity outcomes rather than as first-order engine states. The implementation therefore remains fully compliant with the executable specification while exposing richer diagnostic information to users. :contentReference[oaicite:2]{index=2}

In addition to the scenario suite, the propagation engine is evaluated independently through a declarative harness composed of eleven YAML fixtures. These fixtures validate deterministic fixed-point propagation independently from the Python unit tests. All fixtures converge successfully in one to three iterations, with measured execution times below one millisecond for the synthetic configurations considered. These measurements demonstrate deterministic convergence of the governance algorithm under the evaluated workloads but should not be interpreted as general performance benchmarks. :contentReference[oaicite:3]{index=3}

Overall, the normative evaluation provides executable evidence supporting RQ1 and a substantial part of RQ2. The scenario suite demonstrates that the implemented governance semantics correctly operationalize the principal concepts of the ACM specification, including immutable revisions, baseline management, dependency-aware propagation, runtime replay, controlled dynamic-agent management and governance portability. At the same time, the explicit reporting of implementation deviations illustrates that the evaluation framework is capable of distinguishing specification compliance from implementation-specific modelling decisions, thereby reinforcing the transparency and auditability objectives of the ACM reference model.

\subsection{Cross-Framework Evaluation}
\label{sec:cross_framework_evaluation}

One of the principal objectives of ACM is to provide a framework-independent configuration representation rather than another orchestration framework. Consequently, the evaluation focuses on determining whether heterogeneous native workflows can be normalized into a common ACM representation while preserving their configuration semantics.

The experimental evaluation therefore considers two representative orchestration paradigms. LangGraph exposes an explicit stateful execution graph in which topology, routing and execution structure are directly represented. CrewAI adopts a higher-level orchestration model based on agents, tasks, crews and flows, where part of the execution topology is implicit and resolved by the runtime. These frameworks were deliberately selected because they represent substantially different abstractions while remaining sufficiently mature to evaluate the portability of the ACM reference model. :contentReference[oaicite:1]{index=1}

\begin{table}[t]
\centering
\caption{Characteristics of the evaluated orchestration frameworks.}
\label{tab:framework_characteristics}
\begin{tabularx}{\linewidth}{lXX}
\toprule
 & LangGraph & CrewAI \\
\midrule
Primary abstraction &
Explicit execution graph &
Agents, tasks, crews and flows \\

Execution topology &
Explicit &
Partially implicit \\

Shared execution state &
Native &
Flow state \\

Routing &
Conditional graph edges &
Flow routers and listeners \\

Native runtime replay &
Checkpoint-based &
Flow execution \\

Projection strategy &
Direct extraction &
Flow + Crew reconstruction \\
\bottomrule
\end{tabularx}
\end{table}

The two adapters independently extract native framework objects before projecting them into the ACM intermediate representation described in Section~\ref{sec:projection_pipeline}. Once projected, all validation, propagation and governance algorithms operate exclusively on the normalized ACM representation. Consequently, any difference observed during evaluation reflects the expressive power of the native framework or its introspection capabilities rather than differences in the governance engine itself.

LangGraph provides the most favourable extraction conditions because its workflow topology is explicitly represented. For the evaluated non-trivial workflow, the latest preservation report records complete preservation of nodes, relationships, conditional branches, entry and terminal nodes, prompt references and tool references. Repeated extractions also produce identical canonical digests, demonstrating deterministic projection. The only approximated element concerns the internal semantics of Python routing functions, whose implementation remains intentionally opaque to ACM. :contentReference[oaicite:2]{index=2}

CrewAI requires a different extraction strategy because orchestration semantics are distributed between Flow definitions and Crew objects. The implemented adapter reconstructs a unified representation by combining these complementary sources. According to the latest preservation report, the evaluated Flow+Crew workflow achieves complete structural preservation over the normative evaluation scope, including nodes, relationships, branches, entry points and semantic references. The remaining limitations concern framework internals that are not statically introspectable rather than shortcomings of the ACM metamodel itself. In particular, branch semantics remain represented as abstract routing structures and the internal Flow state schema is currently outside the normative ACM scope. :contentReference[oaicite:3]{index=3}

\begin{table*}[t]
\centering
\caption{Cross-framework projection results on the evaluated workflows.}
\label{tab:projection_results}
\begin{tabular}{lcc}
\toprule
Metric &
LangGraph &
CrewAI Flow+Crew \\
\midrule
Node preservation &
100\% &
100\% \\

Relationship preservation &
100\% &
100\% \\

Branch preservation &
100\% &
100\% \\

Entry point preserved &
Yes &
Yes \\

Terminal nodes preserved &
Yes &
Yes \\

Prompt references &
100\% &
100\% \\

Tool references &
100\% &
100\% \\

Canonical digest stability &
Yes &
Yes \\

Conditional routing semantics &
Approximated &
Approximated \\

State schema &
Supported &
Unsupported \\
\bottomrule
\end{tabular}
\end{table*}

These results provide initial evidence supporting RQ3. Despite relying on substantially different orchestration paradigms, both frameworks can be normalized into a common ACM representation while preserving the structural information required by the governance model. More importantly, ACM explicitly distinguishes between information extracted directly from native framework objects and information supplied by the adapter. This distinction prevents inferred metadata from being presented as native configuration and makes the limits of framework introspection explicit rather than implicit.

The evaluation therefore supports the claim that ACM provides cross-framework portability across two heterogeneous orchestration paradigms. It does not, however, establish universal framework independence. Extending the evaluation to additional orchestration ecosystems remains part of the future work discussed in Section~\ref{sec:future_work}.

\subsection{Information Preservation Analysis}
\label{sec:information_preservation}

Beyond demonstrating that heterogeneous workflows can be projected into a common ACM representation, it is equally important to characterize which aspects of the original configuration are preserved during this transformation. Since different orchestration frameworks expose different levels of introspection, complete semantic equivalence cannot always be expected. Rather than attempting to conceal these differences, ACM explicitly classifies the preservation status of every configuration element.

The preservation analysis is therefore performed at the semantic level rather than the implementation level. Native framework objects are compared with their normalized ACM counterparts and classified according to one of three categories:

\begin{itemize}
    \item \textbf{Preserved}: the native semantics are represented without loss within the ACM metamodel.
    \item \textbf{Approximated}: the information is represented through a normalized ACM construct that preserves its governance semantics while abstracting framework-specific implementation details.
    \item \textbf{Unsupported}: the information falls outside the current normative scope of ACM and is therefore intentionally not represented.
\end{itemize}

This classification differs fundamentally from traditional import fidelity metrics. ACM does not attempt to reproduce every implementation detail of the originating framework. Instead, it preserves those configuration elements that are required for deterministic governance while explicitly documenting any abstraction performed during normalization.

Table~\ref{tab:information_preservation} summarizes the preservation analysis obtained for the evaluated LangGraph and CrewAI workflows.

\begin{table*}[t]
\centering
\caption{Information preservation matrix.}
\label{tab:information_preservation}
\begin{tabular}{lcc}
\toprule
Configuration element &
LangGraph &
CrewAI Flow+Crew \\
\midrule

Workflow topology &
Preserved &
Preserved \\

Workflow nodes &
Preserved &
Preserved \\

Execution relationships &
Preserved &
Preserved \\

Entry point &
Preserved &
Preserved \\

Terminal nodes &
Preserved &
Preserved \\

Prompt references &
Preserved &
Preserved \\

Tool references &
Preserved &
Preserved \\

Model configuration &
Preserved &
Preserved \\

Agent identity &
Preserved &
Preserved \\

Dependency graph &
Preserved &
Preserved \\

Baseline revisions &
Preserved &
Preserved \\

Canonical digests &
Preserved &
Preserved \\

Conditional routing implementation &
Approximated &
Approximated \\

Runtime execution state &
Approximated &
Approximated \\

Internal scheduler metadata &
Unsupported &
Unsupported \\

Framework internal caches &
Unsupported &
Unsupported \\

Python callback implementation &
Unsupported &
Unsupported \\

Native execution engine internals &
Unsupported &
Unsupported \\

\bottomrule
\end{tabular}
\end{table*}

The evaluation shows that all configuration elements required by the ACM governance model are preserved across both evaluated orchestration frameworks. Structural information—including workflow topology, dependency relationships, prompts, tools, models, revisions and configuration identities—is projected without semantic loss.

Only implementation-specific execution mechanisms require abstraction. In particular, conditional routing logic is represented through normalized routing relationships rather than executable Python functions, while runtime execution state is reconstructed from normalized events instead of framework-specific internal objects. These approximations preserve the observable governance semantics without coupling ACM to a particular execution engine.

Finally, several categories of information are intentionally excluded from the ACM representation. Examples include framework scheduler internals, execution caches and implementation-specific callback code. Although these elements influence execution behaviour, they do not constitute governed configuration according to the scope defined by the ACM specification. Their exclusion therefore reflects a deliberate modelling decision rather than a limitation of the projection process.

An additional outcome of this analysis concerns provenance transparency. The projection process distinguishes between information extracted directly from native framework objects and information reconstructed by the projection adapter when explicit native representations are unavailable. Consequently, the provenance of every normalized configuration element remains explicit, allowing auditors to distinguish native configuration from adapter-derived reconstruction. This design improves traceability while avoiding the introduction of implicit or inferred governance metadata.

Overall, the preservation analysis supports RQ3 by demonstrating that ACM provides a stable and framework-independent configuration representation while explicitly documenting the boundaries of its normative scope. Instead of pursuing complete implementation fidelity, ACM preserves the information necessary for deterministic governance and reproducible configuration management, making every approximation and unsupported element explicit.

subsection{Reference Implementation Assessment}
\label{sec:reference_implementation_assessment}

The purpose of the reference implementation is not to introduce a production-ready software platform but to demonstrate that the ACM reference model can be operationalized through a coherent and executable realization. Consequently, this assessment evaluates the correspondence between the normative concepts defined by the specification and their implementation within the Python prototype rather than software quality attributes such as usability, scalability or operational performance.

The assessment considers the principal architectural components introduced in Section~\ref{sec:operationalization}. For each component, we verify that the implementation provides an executable realization of the corresponding normative concept and that this realization is exercised by the automated validation campaign. The objective is therefore to establish implementation completeness with respect to the ACM specification.

\begin{table*}[t]
\centering
\caption{Traceability between specification concepts, implementation components and validation evidence.}
\label{tab:implementation_traceability}
\begin{tabularx}{\textwidth}{l l X}
\toprule
\textbf{Specification concept} &
\textbf{Operational component} &
\textbf{Validation evidence} \\
\midrule

Immutable configuration management &
Revision Manager &
Normative scenarios S01--S04 verify immutable revisions, canonical serialization, deterministic digests and baseline reconstruction. \\

Governance propagation &
Propagation Engine &
Normative scenarios S05--S11 together with eleven declarative propagation fixtures validate deterministic fixed-point convergence, dependency propagation and assurance computation. \\

Runtime governance &
Event Log + Runtime Reducer &
Normative scenarios S12--S17 validate event replay, lifecycle transitions, runtime reconstruction and drift detection. \\

Dynamic agent management &
Runtime Manager &
Normative scenarios S18--S21 validate dynamic agent declaration, permission enforcement and controlled promotion into managed configuration. \\

Cross-framework normalization &
LangGraph Adapter + CrewAI Adapter &
Normative scenarios S22--S27 and the cross-framework evaluation validate deterministic projection, structural preservation and framework-independent governance across both supported orchestration frameworks. \\

\bottomrule
\end{tabularx}
\end{table*}

\subsection{Discussion of Results}
\label{sec:discussion_results}

The evaluation presented in this section provides empirical evidence supporting the feasibility of the proposed Agentic Configuration Management (ACM) reference model. Rather than evaluating a specific orchestration framework, the experiments assess whether the normative concepts introduced by ACM can be implemented, validated and applied consistently across heterogeneous agentic ecosystems.

The reference implementation successfully exercised the complete validation protocol introduced in this work. Twenty-seven normative scenarios verified the principal governance semantics of the specification, including immutable revisions, deterministic propagation, runtime reconstruction, dynamic agent management and framework normalization. These scenario-based evaluations were complemented by declarative propagation fixtures, providing an independent validation of the fixed-point propagation algorithm, and by cross-framework projection experiments performed on representative LangGraph and CrewAI workflows.

The evaluation demonstrates that the ACM metamodel preserves the structural information required for governance independently of the originating orchestration framework. Although certain implementation-specific execution details cannot always be represented directly, the projection process explicitly classifies approximated and unsupported information rather than introducing implicit assumptions. This explicit characterization contributes to the transparency and reproducibility of governance decisions.

The experiments also illustrate the separation between the ACM specification and its reference implementation. The Python prototype should be regarded as one operational realization of the reference model rather than the model itself. The normative concepts remain independent of implementation language, storage technology or execution architecture, suggesting that alternative implementations could remain interoperable provided they preserve the semantics defined by the specification.

The current evaluation nevertheless has several limitations. First, only two orchestration frameworks were evaluated. Although LangGraph and CrewAI represent substantially different orchestration paradigms, they do not cover the full diversity of emerging agentic ecosystems. Second, the experiments focus on semantic correctness rather than operational characteristics such as execution latency, scalability or distributed deployment. Finally, the evaluated workflows remain representative rather than exhaustive. Broader empirical studies involving industrial-scale systems and additional orchestration frameworks remain necessary to further assess the generality of the proposed approach.

Overall, the experimental results support the claim that ACM provides a practical and framework-independent foundation for deterministic configuration governance of agentic systems while clearly identifying the current boundaries of the reference model.

\begin{table*}[t]
\centering
\caption{Summary of the research questions and supporting evidence.}
\label{tab}
\begin{tabularx}{\textwidth}{lXl}
\toprule
\textbf{Research Question} &
\textbf{Evidence obtained} &
\textbf{Assessment} \
\midrule

RQ1 &
The ACM metamodel captures the configuration concepts required for deterministic governance, including immutable revisions, dependency management, governance propagation and runtime reconstruction. These concepts were operationalized by the reference implementation and validated through the normative scenario suite. &
Supported \

RQ2 &
The governance semantics defined by the specification are consistently enforced by the reference implementation. Twenty-seven normative scenarios and eleven propagation fixtures validate deterministic propagation, lifecycle management, assurance computation and drift detection. &
Supported \

RQ3 &
Representative LangGraph and CrewAI workflows were successfully normalized into the ACM representation while preserving the structural information required for governance. Framework-specific execution semantics are explicitly classified as approximated or unsupported when appropriate. &
Supported (evaluated frameworks) \

RQ4 &
The reference implementation demonstrates that the ACM specification can be realized through a coherent software architecture linking immutable configuration management, governance propagation, runtime replay and framework adapters. &
Supported \

\bottomrule
\end{tabularx}
\end{table*}

The results therefore provide consistent evidence supporting the four research questions within the scope of the present evaluation. They should not, however, be interpreted as establishing universal framework independence or production readiness. Instead, they demonstrate that the proposed reference model is operationally feasible, internally consistent and capable of supporting deterministic governance across the evaluated orchestration paradigms. Extending the empirical evaluation to additional frameworks, larger industrial workflows and performance-oriented benchmarks constitutes a natural direction for future work.

\subsection{Threats to Evaluation Validity}
\label{sec:threats_to_validity}

As with any empirical software engineering study, the present evaluation is subject to several limitations. Following established evaluation practice, threats are discussed according to construct, internal, external and conclusion validity.

\subsubsection{Construct Validity}

The evaluation focuses on the governance capabilities defined by the ACM specification rather than on execution performance or orchestration efficiency. Consequently, the selected metrics measure semantic correctness, determinism, information preservation and configuration governance instead of throughput, latency or resource consumption. While these metrics are consistent with the objectives of ACM, they do not provide evidence regarding operational performance.

Furthermore, information preservation is evaluated with respect to the normative ACM metamodel rather than complete implementation fidelity. Certain framework-specific execution details are intentionally classified as approximated or unsupported because they lie outside the scope of governed configuration.

\subsubsection{Internal Validity}

The specification and the reference implementation evolved iteratively throughout the design process, creating a potential risk that implementation decisions could influence the evolution of the reference model. This threat was mitigated by defining the normative concepts independently of implementation details and by validating them through complementary evaluation mechanisms, including scenario-based testing, declarative propagation fixtures and cross-framework projection experiments.

Another potential source of bias arises from the manually developed framework adapters. Errors within an adapter could affect the observed preservation results independently of the ACM metamodel itself. To reduce this risk, the adapters were evaluated on representative workflows exhibiting conditional routing, dependency relationships and non-trivial execution structures.

\subsubsection{External Validity}

The experimental evaluation considers two orchestration frameworks, namely LangGraph and CrewAI. These frameworks were selected because they represent substantially different orchestration paradigms while remaining sufficiently mature to support realistic workflow extraction. Nevertheless, they do not cover the full diversity of emerging agentic ecosystems.

Similarly, the evaluated workflows are representative but remain limited in scale. Although they exercise the principal governance mechanisms of ACM, larger industrial systems may exhibit additional configuration patterns requiring further investigation. Extending the evaluation to additional orchestration frameworks and production-scale deployments therefore constitutes an important direction for future work.

\subsubsection{Conclusion Validity}

The reported results demonstrate that the reference implementation consistently realizes the normative concepts defined by the ACM specification for the evaluated scenarios. However, the evaluation should not be interpreted as demonstrating universal framework independence or production readiness.

Likewise, no statistical claims are made regarding quantitative superiority over existing orchestration frameworks. ACM is proposed as a reference configuration model rather than as an execution framework. Consequently, the empirical evidence supports the operational feasibility and internal consistency of the proposed governance model within the evaluated scope, while broader empirical validation remains necessary to assess its generality.

\section{Discussion}
\label{sec:discusion}
\subsection{From Orchestration to Configuration Governance}

Current agentic frameworks have significantly improved the engineering of AI systems by introducing explicit abstractions for planning, delegation, execution, memory management, and tool orchestration. Frameworks such as LangGraph and CrewAI allow developers to construct increasingly sophisticated workflows while abstracting many low-level implementation concerns. However, despite their architectural differences, these frameworks primarily address \emph{execution semantics}: they specify how autonomous components interact during task execution.

ACM addresses a complementary concern. Rather than defining how agents execute, it defines how agentic systems are represented, versioned, governed, and audited throughout their lifecycle. This distinction is fundamental. While orchestration determines operational behavior at runtime, configuration governance establishes the semantic identity of the system independently of any execution engine.

This separation follows a well-established principle in software engineering. Software Configuration Management (SCM) deliberately distinguishes the managed configuration of a software product from the mechanisms used to build or execute it. Build systems, deployment platforms, and runtime environments evolve independently of SCM because they operate at different abstraction layers. ACM extends this principle to agentic systems by introducing an explicit governance layer above orchestration frameworks.

Figure~\ref{fig:reference_architecture} illustrates this layered organization. Existing frameworks remain responsible for execution, scheduling, communication, and runtime coordination, whereas ACM governs immutable configuration revisions, baselines, provenance, lifecycle semantics, runtime state representation, and governance propagation. Consequently, multiple heterogeneous execution frameworks can share the same governance model without requiring modifications to their internal orchestration mechanisms.

This architectural separation provides several practical advantages.

First, configuration governance becomes framework-independent. Since ACM defines a canonical representation of agentic configurations, governance decisions remain stable even when execution frameworks evolve or are replaced. Migration between orchestration technologies therefore affects execution adapters rather than governance semantics.

Second, lifecycle management becomes deterministic. The state of an agentic configuration is no longer inferred from framework-specific runtime metadata but computed through formally defined lifecycle rules and propagation semantics. This guarantees that governance decisions remain reproducible independently of execution histories.

Third, provenance information becomes explicit. Existing orchestration frameworks generally expose execution traces, logs, or runtime events. Although these artifacts provide valuable operational observability, they do not necessarily preserve the complete evolution of configuration objects or the rationale behind governance decisions. ACM complements runtime tracing by maintaining immutable configuration revisions together with explicit relationships between design-time artifacts and runtime instances.

Finally, separating governance from orchestration improves interoperability. Since governance semantics are defined independently of execution semantics, heterogeneous frameworks can be projected into a common representation while preserving their execution-specific capabilities. This separation is analogous to the role played by intermediate representations in compiler design or by platform-independent configuration models in cloud-native infrastructures.

The objective of ACM is therefore not to replace existing orchestration frameworks, nor to compete with them. Instead, ACM provides a governance abstraction that operates orthogonally to execution. Execution frameworks remain responsible for computation, scheduling, and coordination, whereas ACM establishes the formal configuration model required for version control, reproducibility, traceability, compliance, and long-term lifecycle management.

This distinction also clarifies the scientific contribution of ACM. Existing frameworks primarily optimize \emph{how} autonomous systems execute, whereas ACM formalizes \emph{what} constitutes the governed configuration of an agentic system and \emph{how} that configuration evolves over time under explicit governance rules. The contribution therefore resides at the level of configuration semantics rather than execution semantics, making ACM complementary to existing orchestration technologies rather than an alternative to them.

\subsection{Implications for AI Governance}

As agentic systems become increasingly autonomous, software engineering concerns extend beyond execution correctness toward broader governance requirements. Organizations deploying AI systems must not only ensure functional behavior but also demonstrate how configurations evolve, how decisions can be reproduced, and how responsibilities can be traced throughout the system lifecycle. These requirements are commonly associated with AI governance frameworks, which emphasize accountability, transparency, auditability, and reproducibility as fundamental properties of trustworthy AI systems.

ACM contributes to these objectives by introducing governance mechanisms at the configuration level rather than at the execution level. Instead of attempting to regulate agent reasoning or model inference, ACM governs the evolution of the artifacts that define an agentic system. This distinction is important because configuration artifacts remain observable, versionable, and verifiable independently of the non-deterministic behavior of large language models.

A first contribution concerns \emph{accountability}. Every governed object in ACM evolves through immutable revisions linked by explicit lineage relationships. Configuration changes therefore become attributable to identifiable revision events rather than implicit modifications of mutable objects. Although ACM does not assign organizational responsibilities or access-control policies, it provides the structural foundation required to associate governance decisions with concrete configuration states.

A second contribution concerns \emph{traceability}. Modern agentic frameworks typically expose runtime traces describing execution flows, tool invocations, or message exchanges. These traces explain what occurred during execution but generally do not preserve the complete configuration history that produced those executions. ACM complements runtime observability by maintaining explicit relationships between configuration revisions, baselines, projected runtime instances, and governance states. Consequently, execution artifacts can be interpreted within the context of the configuration from which they originated.

Closely related is the notion of \emph{auditability}. Since governance decisions are computed from immutable revisions and formally defined propagation semantics, the complete governance state of a managed system can be reconstructed from its recorded configuration history. This property enables independent verification of lifecycle states, dependency relationships, and governance decisions without requiring access to the original execution environment. Auditability therefore emerges as a consequence of deterministic governance semantics rather than from additional monitoring mechanisms.

Reproducibility constitutes another important implication. The stochastic nature of LLM-based systems makes behavioral reproducibility inherently difficult, as identical prompts may lead to different outputs under identical configurations. ACM does not attempt to eliminate this source of variability. Instead, it guarantees the reproducibility of the governed configuration itself. Given the same immutable revisions, dependency graph, lifecycle definitions, and governance rules, the resulting governance state is deterministic. This distinction separates configuration reproducibility from behavioral reproducibility, two concepts that are frequently conflated in current discussions on agentic AI.

The proposed governance model also improves long-term maintainability. Because configuration semantics are decoupled from execution semantics, governance information remains stable even when orchestration frameworks, runtime implementations, or underlying language models evolve. This separation allows organizations to preserve the historical evolution of their agentic systems despite technological changes affecting the execution layer.

It is important to clarify the scope of these contributions. ACM is not intended to provide a complete AI governance framework. It does not address organizational governance processes, ethical decision-making, regulatory compliance, risk assessment, model evaluation, security policies, or human oversight mechanisms. Instead, ACM contributes a formal configuration governance layer that can support such governance activities by providing immutable configuration histories, explicit provenance relationships, deterministic lifecycle semantics, and reproducible governance decisions.

From this perspective, ACM should be viewed as a foundational infrastructure for AI governance rather than as a governance methodology. Just as Software Configuration Management became an essential component of software engineering without replacing software development processes, configuration governance may become a foundational capability for future agentic systems while remaining complementary to higher-level governance frameworks, organizational policies, and regulatory requirements.

\subsection{Interoperability and Standardization}

One of the principal motivations behind ACM is to reduce the fragmentation currently observed across agentic software ecosystems. Contemporary orchestration frameworks provide different execution models, configuration formats, lifecycle abstractions, and runtime representations. While this diversity has accelerated innovation, it also complicates the portability, governance, and long-term management of agentic systems. Configuration artifacts developed for one framework generally cannot be interpreted directly by another without framework-specific transformations.

ACM approaches this problem by introducing a framework-independent configuration model that is intentionally decoupled from execution semantics. Rather than defining how workflows should be executed, ACM specifies a canonical representation of the governed entities composing an agentic system, together with their relationships, lifecycle semantics, provenance, and governance state. Execution-specific concepts remain encapsulated within projection adapters, while governance semantics remain stable across heterogeneous orchestration technologies.

This separation establishes a common semantic vocabulary for describing agentic configurations. Core concepts such as configuration items, immutable revisions, baselines, dependency relationships, runtime instances, lifecycle states, and governance propagation become independent of the implementation choices made by individual execution frameworks. Consequently, heterogeneous systems can be represented using equivalent governance semantics despite relying on different orchestration paradigms.

The evaluation presented in Section~7 provides an initial illustration of this principle. Both LangGraph and CrewAI, despite exposing distinct execution models, can be projected onto the same ACM representation while preserving the governance properties defined by the reference model. This result does not imply semantic equivalence between the frameworks themselves; rather, it demonstrates that governance semantics can be abstracted from execution semantics through a common intermediate representation.

The proposed architecture therefore resembles the role played by intermediate representations in other engineering disciplines. In compiler design, multiple programming languages can target a common intermediate representation before code generation. Similarly, infrastructure-as-code technologies frequently rely on declarative resource models that remain independent of specific cloud providers. ACM applies an analogous principle to agentic systems by separating governance representation from execution implementation.

Such a representation also creates opportunities for interoperability beyond framework migration. Because governance information is expressed independently of execution technologies, configuration repositories, governance dashboards, compliance tools, provenance systems, and auditing platforms could exchange governed configurations without requiring direct knowledge of the underlying orchestration framework. This interoperability is particularly valuable for organizations operating heterogeneous AI ecosystems where multiple execution technologies coexist.

These characteristics suggest that ACM may provide the foundations of a common configuration language for governed agentic systems. Nevertheless, the objective of this work is not to propose a standard. Standardization requires broad industrial adoption, interoperability across a significantly larger ecosystem of frameworks, and community consensus extending beyond the scope of a single research contribution. The present work should therefore be understood as a candidate reference model whose formal semantics are intended to facilitate future discussions on interoperability rather than prescribe a universal representation.

Future validation of this hypothesis will require evaluating ACM against additional orchestration frameworks, execution paradigms, and deployment environments. Such studies would determine whether the abstractions introduced by ACM are sufficiently expressive to capture the diversity of emerging agentic architectures while preserving the framework independence that constitutes the primary objective of the proposed governance model.

\subsection{Scope and Current Limitations}

The objective of ACM is deliberately focused. Rather than proposing a complete framework for building, deploying, or operating agentic systems, ACM defines a formal configuration governance model that specifies how agentic configurations are represented, versioned, governed, and related throughout their lifecycle. This scope intentionally limits the contribution to the governance layer while remaining independent of execution technologies.

Consequently, ACM does not attempt to replace existing orchestration frameworks. Components responsible for task scheduling, agent communication, planning strategies, tool execution, memory management, distributed coordination, or inference remain under the responsibility of execution engines such as LangGraph or CrewAI. ACM neither prescribes nor constrains the internal execution semantics of these frameworks; instead, it provides a common governance representation that can coexist with heterogeneous execution models.

Similarly, ACM does not address the internal behavior of foundation models. The proposed model is independent of the reasoning capabilities, prompting strategies, or inference algorithms employed by large language models. Questions related to model alignment, hallucination mitigation, prompt optimization, retrieval quality, or reasoning performance remain outside the scope of the present work. These aspects influence the behavior of agentic systems but do not modify the governance semantics of their configurations.

The proposed model should also not be interpreted as a complete AI governance framework. ACM does not formalize organizational policies, human oversight processes, ethical principles, regulatory compliance mechanisms, security controls, access management, or risk assessment methodologies. Such concerns require organizational and domain-specific governance processes extending beyond configuration management. Instead, ACM provides foundational governance primitives—immutable revisions, provenance relationships, lifecycle semantics, and deterministic propagation—that can support higher-level governance frameworks without replacing them.

The current evaluation also imposes several limitations on the generalizability of the results. Experimental validation has been conducted using two representative orchestration frameworks, LangGraph and CrewAI, selected because they expose distinct orchestration paradigms while remaining sufficiently mature for reproducible experimentation. Although these results demonstrate the feasibility of framework-independent governance across heterogeneous execution models, they do not establish universal interoperability across the rapidly evolving ecosystem of agentic platforms.

Likewise, the reference implementation should be viewed as a proof of implementability rather than as a production-ready governance platform. Its primary purpose is to demonstrate that the formal concepts introduced by ACM can be operationalized consistently and evaluated experimentally. Questions related to large-scale distributed deployments, high-availability architectures, enterprise integration, and long-term operational optimization remain beyond the scope of the current implementation.

Finally, the evolution of agentic AI represents an additional source of uncertainty. New orchestration paradigms, autonomous coordination mechanisms, dynamic agent creation strategies, and execution models continue to emerge at a rapid pace. Although ACM has been designed around abstractions intended to remain independent of specific frameworks, additional validation will be necessary to determine whether these abstractions remain sufficiently expressive as the agentic ecosystem evolves.

These limitations are intentional rather than accidental. By restricting its scope to configuration governance, ACM aims to establish a stable and formally defined foundation upon which richer governance capabilities, broader interoperability studies, and future standardization efforts may progressively be built. This design choice favors conceptual clarity and extensibility over exhaustive coverage of every aspect of agentic AI systems.

\subsection{Research Perspectives}

The rapid evolution of agentic AI suggests that configuration governance will become an increasingly important research area as autonomous systems continue to grow in complexity, scale, and autonomy. While ACM establishes a formal foundation for representing and governing agentic configurations, it also raises several research questions extending beyond the scope of the present work.

A first research direction concerns the formal characterization of governance semantics. The lifecycle model, dependency propagation rules, and state transition mechanisms introduced by ACM constitute an initial semantic framework. Future research could investigate their formal properties, including soundness, completeness, convergence, consistency preservation, and compositionality. Such analyses would strengthen the theoretical foundations of configuration governance and facilitate formal verification of governance algorithms.

A second perspective concerns interoperability across heterogeneous ecosystems. The experiments presented in this work demonstrate the feasibility of projecting two distinct orchestration frameworks onto a common governance model. A broader research challenge consists of determining whether a minimal set of governance abstractions can remain expressive across the diversity of emerging agentic architectures, including decentralized multi-agent systems, dynamically generated workflows, hierarchical orchestration models, and future execution paradigms. Addressing this question would contribute to understanding the limits of framework-independent governance representations.

Another promising direction concerns the relationship between design-time configurations and runtime evolution. Agentic systems increasingly exhibit adaptive behaviors, dynamic agent creation, autonomous workflow restructuring, and continuous interaction with external environments. Understanding how governance models should represent and preserve provenance under these evolving conditions remains an open research problem. In particular, future work may investigate whether runtime adaptation itself should become a first-class governed configuration artifact rather than merely an execution event.

Configuration governance also opens new perspectives for empirical software engineering. The availability of immutable configuration histories, explicit dependency graphs, and deterministic lifecycle semantics creates opportunities for quantitative studies of configuration evolution, governance complexity, configuration drift, and maintainability of agentic systems. Such empirical investigations could contribute to establishing evidence-based engineering practices comparable to those that emerged in traditional Software Configuration Management.

Beyond software engineering, ACM may provide a foundation for interdisciplinary research on trustworthy autonomous systems. Configuration governance naturally intersects with provenance, accountability, auditability, digital forensics, compliance engineering, and AI assurance. Although ACM intentionally focuses on configuration semantics, its formal representation could serve as a common substrate supporting research across these complementary domains.

Finally, the emergence of governance-oriented reference models raises broader questions regarding the future engineering of agentic systems. As orchestration frameworks continue to diversify and autonomous capabilities expand, execution technologies may evolve independently from governance infrastructures. Investigating this architectural separation, and determining the extent to which configuration governance should become an independent engineering discipline analogous to Software Configuration Management, represents a long-term research question motivated by the present work.

Rather than presenting configuration governance as a completed solution, ACM should therefore be viewed as an initial step toward a broader scientific agenda. By introducing a formal representation of governed agentic configurations, the proposed reference model aims to establish a conceptual foundation upon which future theoretical developments, empirical studies, interoperability initiatives, and governance methodologies can progressively be built.

\section{Validity of the Proposed Model and Evaluation}
\subsection{Validity of the Reference Model}

The primary objective of ACM is not to provide an exhaustive representation of every existing or future agentic system, but rather to define a minimal and framework-independent reference model for configuration governance. Consequently, the validity of the proposed model depends on whether its abstractions are sufficiently expressive to capture the essential governance concepts shared across heterogeneous agentic ecosystems while remaining independent of implementation-specific execution semantics.

The reference model was developed by extending well-established principles from Software Configuration Management to the emerging domain of agentic AI. Core abstractions—including configuration items, immutable revisions, baselines, dependency relationships, lifecycle states, runtime instances, and governance propagation—were selected because they represent concepts that can be identified independently of any particular orchestration framework. Rather than reproducing framework-specific execution constructs, ACM deliberately models only those entities required to govern the evolution of agentic configurations throughout their lifecycle.

The evaluation presented in this work provides initial evidence supporting the adequacy of these abstractions. Despite their different orchestration paradigms, both LangGraph and CrewAI can be projected onto the same governance representation without modifying the underlying reference model. This result suggests that the proposed abstractions capture governance concepts that are not inherently tied to a particular execution technology. At the same time, the experiments should not be interpreted as demonstrating universality, since only a limited subset of contemporary frameworks has been investigated.

The proposed model also intentionally favors conceptual simplicity over exhaustive expressiveness. Several emerging aspects of agentic systems, including autonomous agent generation, decentralized coordination, adaptive execution strategies, and continuously evolving runtime topologies, are represented only to the extent necessary to preserve governance semantics. More specialized execution concepts remain delegated to framework-specific adapters rather than incorporated directly into the reference model. This design choice reduces conceptual complexity while preserving extensibility through future model evolution.

Another important aspect of validity concerns the separation between governance semantics and execution semantics. By restricting the reference model to configuration governance, ACM avoids embedding assumptions regarding planning algorithms, scheduling policies, communication protocols, reasoning mechanisms, or language-model behavior. This separation reduces coupling with rapidly evolving execution technologies and contributes to the long-term stability of the governance abstractions. Nevertheless, additional empirical validation will be necessary to determine whether this separation remains appropriate as new categories of agentic architectures emerge.

Finally, the proposed reference model should be regarded as an evolving conceptual framework rather than as a definitive ontology for agentic systems. The abstractions introduced in ACM represent a first attempt to formalize configuration governance independently of orchestration technologies. Their validity therefore rests not on claims of completeness, but on their ability to provide a coherent, extensible, and operationally meaningful representation of governed agentic configurations. Future validation across a broader range of execution frameworks, application domains, and deployment environments will ultimately determine the generality of the proposed reference model.

\subsection{Validity of the Formal Semantics}

The formal semantics introduced in Section~5 aim to provide a rigorous and deterministic interpretation of the governance concepts defined by the ACM reference model. Rather than serving as an executable specification of agent execution, these formalizations describe the evolution of governed configurations through explicit mathematical structures governing lifecycle transitions, dependency propagation, runtime representation, and governance state computation. Their validity therefore depends on whether they constitute a coherent and faithful formalization of the conceptual model presented in the preceding sections.

To improve consistency, the formal semantics were developed from a common set of mathematical primitives. Configuration entities are represented as immutable objects associated with explicit revisions, while lifecycle transitions, dependency relations, and governance propagation are defined through deterministic functions over these objects. This unified formalization reduces ambiguities that frequently arise when lifecycle management, provenance, and dependency reasoning are described only informally.

An important property of the proposed semantics is determinism. Given an identical configuration graph, the same revision history, and the same governance rules, the resulting governance state is uniquely determined. Consequently, governance decisions do not depend on implementation-specific behaviors or execution environments but only on the formally defined state of the governed configuration. This determinism is essential for reproducibility, auditability, and consistent lifecycle management.

The formalization also preserves a clear separation between governance semantics and execution semantics. The mathematical model intentionally abstracts away execution-specific mechanisms such as scheduling policies, communication protocols, planning strategies, or language-model inference. Instead, runtime information is represented only insofar as it contributes to configuration governance. This restriction allows the formal semantics to remain independent of particular orchestration technologies while maintaining compatibility with heterogeneous execution models.

Nevertheless, the current formalization deliberately remains conservative. The proposed semantics describe deterministic governance processes and therefore do not explicitly model uncertainty, probabilistic reasoning, temporal constraints, or stochastic execution behaviors that naturally arise in LLM-based systems. Likewise, the formal model assumes that runtime observations are correctly projected into the governance layer by framework-specific adapters. The correctness of these projections is evaluated experimentally but is not formally verified within the semantic model itself.

Another limitation concerns formal verification. Although the mathematical definitions are internally consistent and have been operationalized through the reference implementation, the present work does not provide formal proofs of properties such as soundness, completeness, minimality, termination, or semantic equivalence between projected configurations and their original framework representations. Establishing such properties would require dedicated theoretical analyses extending beyond the objectives of this paper.

The proposed formal semantics should therefore be regarded as an initial mathematical foundation for configuration governance rather than as a complete formal theory of agentic systems. Their primary contribution is to provide an unambiguous semantic interpretation of the ACM governance model while remaining sufficiently simple to support practical implementation and experimental evaluation. Future work may extend these semantics through formal verification, richer logical frameworks, temporal reasoning, and additional theoretical analyses to further strengthen the mathematical foundations of governed agentic systems.

\subsection{Generalizability}

The experimental evaluation presented in Section~7 provides an initial validation of the ACM reference model through representative governance scenarios implemented on two heterogeneous orchestration frameworks. As discussed in the evaluation, the obtained results demonstrate the feasibility of projecting framework-specific configurations onto a common governance representation and validating the proposed lifecycle, provenance, and propagation semantics. Nevertheless, these experiments should be interpreted as evidence of feasibility rather than proof of universal applicability.

The generalizability of the proposed model primarily depends on the stability of its underlying abstractions rather than on the specific execution frameworks considered in this work. ACM intentionally defines governance concepts that are independent of orchestration mechanisms, including immutable configuration revisions, dependency relationships, lifecycle semantics, governance propagation, and runtime configuration representations. Since these abstractions are not derived from framework-specific execution primitives, they are expected to remain applicable across heterogeneous agentic technologies provided that suitable projection mechanisms can be established.

At the same time, the current evaluation covers only a limited portion of the rapidly evolving agentic ecosystem. LangGraph and CrewAI were selected because they represent two widely adopted orchestration paradigms with sufficiently different architectural characteristics to evaluate framework-independent governance. However, numerous execution models remain outside the scope of the present work, including decentralized multi-agent platforms, event-driven architectures, dynamically synthesized workflows, and future orchestration paradigms that may introduce governance requirements not yet represented within ACM.

Generalizability must also be considered with respect to application domains. The proposed reference model was designed to describe governance properties independently of any particular business context. Consequently, the semantics introduced by ACM are expected to remain applicable regardless of whether agentic systems are deployed for software engineering, knowledge management, scientific workflows, industrial automation, customer support, or other application domains. Nevertheless, domain-specific governance constraints may require extensions to the current model while preserving its underlying configuration semantics.

Another important aspect concerns scalability. Although the governance algorithms defined by ACM are independent of graph size, the experimental evaluation remains limited to representative workflows of practical complexity. Additional empirical studies involving substantially larger configuration graphs, more complex dependency structures, and long-running agentic systems will be necessary to characterize the scalability of the proposed governance model under industrial operating conditions.

Finally, the rapid evolution of agentic AI represents an inherent challenge for any reference model. New execution paradigms, orchestration mechanisms, and autonomous coordination strategies continue to emerge at a pace that exceeds traditional software engineering cycles. The long-term generalizability of ACM will therefore depend on whether its governance abstractions remain sufficiently stable to accommodate these technological evolutions without requiring fundamental modifications to the reference model.

For these reasons, the conclusions presented in this work should be interpreted within the scope of the evaluated frameworks and scenarios. The experimental results provide encouraging evidence that configuration governance can be defined independently of execution technologies, but broader validation across additional frameworks, application domains, and deployment environments will ultimately determine the practical generality of the proposed reference model.

\subsection{Reproducibility and Replicability}

Reproducibility and replicability constitute essential requirements for scientific research as well as for the engineering of governed software systems. Although these terms are sometimes used interchangeably, they refer to distinct validation objectives. Reproducibility concerns the ability to obtain identical results using the same artifacts, implementation, experimental protocol, and execution environment, whereas replicability evaluates whether comparable conclusions can be reached through an independent implementation following the same conceptual specification.

The ACM reference model has been designed to facilitate reproducibility at the governance level. Configuration revisions are immutable, governance rules are formally specified, lifecycle transitions are deterministic, and dependency propagation follows explicitly defined semantics. Consequently, given the same governed configuration, revision history, and governance rules, the resulting governance state is uniquely determined. This property enables independent verification of governance decisions without ambiguity arising from mutable configuration artifacts.

It is important, however, to distinguish governance reproducibility from behavioral reproducibility. Agentic systems frequently rely on large language models whose inference processes are inherently stochastic. Variability introduced by probabilistic decoding, model updates, external services, or runtime environments may lead to different execution outputs despite identical configurations. ACM does not attempt to eliminate this source of variability. Instead, it guarantees that the governed configuration from which these executions originate remains stable, traceable, and reproducible. In this sense, ACM provides reproducibility of the managed configuration rather than reproducibility of the autonomous behavior itself.

The reference implementation and experimental protocol further contribute to reproducibility by providing a formal specification of the governance model, deterministic evaluation scenarios, explicit lifecycle definitions, and reproducible projection procedures for the evaluated frameworks. These elements enable other researchers to reproduce the reported governance results under equivalent experimental conditions while independently verifying the consistency of the proposed semantics.

Replicability represents a stronger validation criterion. Although the reference implementation demonstrates that the ACM model can be operationalized successfully, the present work does not establish that identical conclusions would necessarily be obtained through an entirely independent implementation. Such validation would require separate implementations developed by different research teams, potentially targeting additional orchestration frameworks and execution environments while adhering to the same formal specification.

From this perspective, the mathematical formalization introduced in Section~5 plays a central role. By defining governance semantics independently of implementation details, the formal specification provides a stable reference against which independent implementations can be evaluated. The validity of the ACM model therefore does not rely exclusively on the correctness of the current software prototype but also on the ability of future implementations to realize the same formal semantics and reach equivalent governance conclusions.

Overall, ACM establishes the conditions necessary for reproducible configuration governance while providing a formal foundation that supports future replication studies. Demonstrating replicability across independent implementations, additional frameworks, and diverse execution environments remains an important direction for future validation of the proposed reference model.

\subsection{Mitigation Strategies}

Although the limitations discussed in the previous sections cannot be entirely eliminated, several methodological and architectural decisions were deliberately adopted to reduce their impact on the validity of the proposed model and its experimental evaluation.

A first mitigation strategy consists in the explicit separation between the conceptual model, its formal semantics, and its software implementation. The ACM reference model was defined independently of any particular orchestration framework, while the mathematical formalization introduced in Section~5 specifies its governance semantics without relying on implementation-specific behaviors. The reference implementation therefore represents an operational realization of the model rather than the model itself. This separation reduces the risk that implementation artifacts influence the conceptual validity of the proposed governance abstractions.

A second strategy concerns the use of deterministic governance semantics. Lifecycle transitions, dependency propagation, provenance relationships, and governance state computation are defined through explicit mathematical rules operating on immutable configuration revisions. Consequently, governance decisions remain independent of execution order, runtime environments, or implementation-specific optimizations, thereby reducing ambiguity during experimental validation and facilitating independent verification of governance results.

The experimental protocol also contributes to mitigating validity threats. Evaluation scenarios were derived from the governance properties defined by the reference model rather than being designed around framework-specific capabilities. This approach reduces the likelihood of overfitting the evaluation to a particular execution technology and ensures that the experiments remain aligned with the conceptual objectives of ACM. Furthermore, the evaluation was conducted on two heterogeneous orchestration frameworks exhibiting different execution paradigms, providing initial evidence that the proposed governance abstractions are not inherently coupled to a single implementation ecosystem.

Another important mitigation strategy is the adoption of immutable configuration revisions as the foundation of governance. By preventing in-place modification of governed artifacts, ACM ensures that provenance relationships, lifecycle transitions, and governance decisions remain traceable throughout the complete evolution of an agentic system. This immutability simplifies independent auditing and reduces the possibility of inconsistencies introduced by mutable configuration states.

The publication of a formal specification together with a reference implementation further contributes to scientific transparency. The mathematical definitions, governance algorithms, and evaluation protocol collectively provide a reproducible basis for independent verification and future replication studies. Although additional implementations and broader experimental campaigns remain necessary to strengthen external validity, the current work establishes an explicit reference against which future developments can be objectively compared.

Finally, the deliberate restriction of ACM to configuration governance constitutes a methodological mitigation rather than a limitation alone. By clearly separating governance semantics from execution semantics, AI reasoning, organizational governance, and regulatory compliance, the proposed model avoids introducing assumptions that cannot be formally justified within the scope of the present work. This focused design reduces conceptual ambiguity while providing a stable foundation upon which richer governance capabilities may subsequently be developed.

Taken together, these methodological choices do not eliminate the validity threats identified throughout this section, but they substantially reduce their impact on the interpretation of the reported results. More importantly, they establish a transparent scientific framework that enables the proposed reference model to be critically evaluated, independently implemented, and progressively refined as the field of agentic AI continues to evolve.

\section{Future Work}
\label{sec:future_work}

\subsection{Extension of the Reference Model}

The ACM reference model has been intentionally designed around a minimal set of governance abstractions that capture the fundamental concepts required to represent, version, and govern agentic configurations. This deliberate minimality favors conceptual clarity and framework independence while providing a stable foundation for future extensions. As the field of agentic AI continues to evolve, the reference model is expected to progressively incorporate additional governance concepts without modifying its core architectural principles.

One natural direction concerns the enrichment of governed configuration entities. The current model focuses on configuration items, immutable revisions, baselines, dependency relationships, runtime representations, and governance states. Future versions of ACM may introduce additional first-class governance objects describing richer operational artifacts, long-lived execution contexts, shared resources, collaborative knowledge structures, or other configuration entities emerging from future orchestration paradigms. Such extensions should preserve the separation between governance semantics and execution semantics established by the present work.

A second direction involves extending governance relationships. While the current reference model captures structural dependencies and lifecycle propagation, future versions may introduce richer semantic relationships expressing compatibility constraints, policy inheritance, conditional dependencies, assurance requirements, or domain-specific governance rules. Expanding the relationship model would enable more expressive governance reasoning while maintaining the deterministic principles that characterize ACM.

Another important perspective concerns increasingly dynamic agentic systems. Current orchestration frameworks are beginning to support runtime adaptation mechanisms, including autonomous agent creation, dynamic workflow restructuring, and self-organizing execution graphs. Although ACM already provides governance mechanisms for runtime representations, future versions of the reference model may explicitly incorporate governance abstractions dedicated to dynamically evolving configurations while preserving immutable configuration histories and deterministic governance semantics.

The evolution of governance policies also represents a promising area for future development. The present work primarily formalizes lifecycle management, provenance, dependency propagation, and runtime configuration semantics. Future extensions could support richer policy languages capable of expressing organizational governance constraints, configurable assurance requirements, policy composition mechanisms, or domain-specific governance rules while remaining compatible with the framework-independent philosophy of ACM.

Future versions of the reference model may also benefit from additional semantic dimensions describing governed configurations. Beyond lifecycle and runtime semantics, future extensions could introduce explicit representations of governance maturity, configuration quality, validation status, assurance evidence, or other metadata supporting advanced governance analyses. Such additions should remain orthogonal to execution semantics to preserve the modular structure of the proposed reference model.

Throughout these extensions, maintaining backward compatibility of the core governance abstractions will remain an important design objective. Rather than continuously redefining the conceptual foundations of ACM, future evolution should favor incremental enrichment of the reference model through modular extensions that preserve the stability of its fundamental semantics. This approach is consistent with the historical evolution of Software Configuration Management, where stable conceptual foundations have supported decades of methodological and technological innovation.

Overall, the future evolution of ACM is expected to emphasize extensibility rather than conceptual replacement. By preserving a minimal governance core while progressively integrating richer configuration abstractions, the reference model can evolve alongside emerging agentic technologies without compromising the framework independence and governance semantics that constitute its primary contribution.

\subsection{Broader Framework Validation}

The experimental validation presented in this work demonstrates the feasibility of applying ACM to two representative orchestration frameworks exhibiting distinct execution paradigms. Although these results provide encouraging evidence of framework-independent configuration governance, they represent only an initial validation of the proposed reference model. A broader evaluation across the rapidly expanding agentic ecosystem will be required to assess the expressiveness, robustness, and long-term applicability of the ACM abstractions.

A first direction consists of extending the validation to additional orchestration frameworks. Emerging platforms such as OpenAI Agents SDK, Microsoft AutoGen, Semantic Kernel, or future agent orchestration environments expose different execution models, communication mechanisms, and lifecycle abstractions. Evaluating ACM against these heterogeneous systems would provide valuable insight into the stability of the proposed governance model and help identify abstractions requiring further refinement. Such studies would also enable quantitative assessment of projection fidelity, semantic preservation, and potential information loss across increasingly diverse execution paradigms.

Beyond orchestration frameworks, future validation should also consider the broader AgentOps ecosystem. Modern agentic applications increasingly rely on complementary platforms dedicated to observability, tracing, evaluation, telemetry, prompt management, experimentation, and operational monitoring. Systems such as Helicone, LangSmith, OpenTelemetry-based observability pipelines, or ML lifecycle platforms provide metadata describing runtime behavior, execution traces, evaluation results, cost metrics, and operational characteristics that are highly relevant for governance activities.

Although these platforms do not directly define agentic configurations, they produce governance-relevant information that complements the configuration model introduced by ACM. Future work should therefore investigate how operational evidence generated by such systems could be associated with immutable configuration revisions, runtime instances, and governance states without violating the separation between execution semantics and configuration governance established by the reference model.

Another important validation perspective concerns interoperability between governance and operational tooling. Rather than replacing existing AgentOps platforms, ACM could act as a common governance layer capable of linking heterogeneous sources of operational evidence to a stable configuration representation. Such an approach would facilitate end-to-end provenance, enabling organizations to relate configuration evolution, execution traces, evaluation metrics, and governance decisions within a unified conceptual model while preserving the specialized role of each operational platform.

Future validation should also extend beyond laboratory-scale experiments toward industrial deployment scenarios. Long-running autonomous systems, collaborative multi-agent environments, continuously evolving workflows, and enterprise-scale governance infrastructures represent important validation contexts that cannot be fully addressed through the current experimental campaign. Evaluating ACM under these operational conditions would provide valuable evidence regarding scalability, maintainability, governance consistency, and operational integration.

Ultimately, broader validation is expected to assess not only the compatibility of ACM with additional orchestration frameworks, but also its ability to serve as a governance abstraction spanning the wider AgentOps ecosystem. Demonstrating that a single governance model can consistently integrate configuration management, provenance, runtime representations, observability metadata, and operational evidence across heterogeneous technologies would substantially strengthen the practical relevance of the proposed reference model.

\subsection{Formal Verification of Governance Semantics}

The mathematical formalization introduced in Section~5 provides a precise semantic foundation for the ACM reference model. However, the present work primarily demonstrates the operational consistency of these semantics through their implementation and experimental evaluation. An important direction for future work is the formal verification of the governance properties defined by the model.

Several fundamental properties of the proposed governance semantics could be investigated through formal methods. For example, lifecycle state machines may be analyzed to verify the absence of invalid transition sequences, while governance propagation algorithms could be studied to establish properties such as determinism, convergence, termination, and consistency preservation. Similarly, the dependency model could be analyzed to ensure that governance decisions remain invariant under equivalent configuration representations and that propagation produces unique governance states for identical configuration graphs.

Another important perspective concerns the verification of invariants associated with immutable configuration management. The ACM reference model assumes that revisions remain immutable, provenance relationships remain consistent, and governance decisions preserve traceability throughout the complete lifecycle of governed configurations. Formal verification techniques could provide mathematical guarantees that these invariants cannot be violated by compliant implementations of the governance model.

Future work may also investigate the semantic equivalence between framework-specific configurations and their ACM representations. Projection adapters currently translate heterogeneous orchestration artifacts into a common governance model while preserving their governance semantics. Establishing formal equivalence criteria between source configurations and projected representations would strengthen confidence in the correctness of these transformations and provide a rigorous basis for measuring semantic preservation and information loss.

Another promising direction concerns the specification of governance properties using formal specification languages and verification techniques. State-transition systems, temporal properties, graph invariants, and declarative constraints could be formally expressed and automatically verified against the governance model. Such approaches would facilitate independent validation of ACM implementations while reducing ambiguity in the interpretation of governance semantics.

Beyond individual properties, formal verification may contribute to establishing stronger theoretical guarantees regarding the ACM reference model itself. Questions related to semantic soundness, completeness with respect to the intended governance model, compositionality of governance rules, and correctness of propagation mechanisms remain open research problems. Addressing these questions would significantly strengthen the theoretical foundations of configuration governance for agentic systems.

The progressive integration of formal verification into ACM would also benefit practical adoption. Verified governance semantics could increase confidence in implementations deployed within environments requiring strong guarantees regarding traceability, configuration integrity, reproducibility, and auditability. While such guarantees cannot eliminate uncertainties associated with autonomous AI behavior, they can substantially improve confidence in the correctness of the configuration governance layer supporting these systems.

Ultimately, the mathematical formalization proposed in this work should be viewed as the first step toward a verifiable theory of configuration governance for agentic systems. Future formal analyses are expected to complement the experimental validation presented in this paper by providing mathematical guarantees that further reinforce the robustness and long-term reliability of the ACM reference model.

\subsection{Toward an Agentic Configuration Ecosystem}

The long-term value of a reference model extends beyond its conceptual specification. If adopted by a sufficiently broad ecosystem of tools and platforms, a common governance model can facilitate interoperability, reduce integration complexity, and encourage the emergence of reusable governance services. In this perspective, ACM should not be viewed solely as an isolated configuration model, but as a potential foundation for a broader ecosystem dedicated to the governance of agentic systems.

A first opportunity concerns the development of governance-aware tooling. Configuration editors, visualization environments, validation engines, lifecycle managers, provenance explorers, and governance dashboards could operate directly on ACM representations while remaining independent of any specific orchestration framework. Such tools would provide a consistent user experience across heterogeneous execution technologies by relying on a common governance vocabulary rather than framework-specific configuration formats.

Another perspective involves the integration of ACM with existing AgentOps infrastructures. Contemporary operational platforms increasingly provide execution tracing, observability, prompt management, telemetry, evaluation pipelines, and cost monitoring. Rather than duplicating these capabilities, ACM could provide the stable governance layer linking operational evidence to immutable configuration revisions. This separation would preserve the specialized responsibilities of operational platforms while enabling configuration-centric reasoning across the complete lifecycle of agentic systems.

The establishment of common governance repositories also represents a promising direction. Instead of maintaining framework-specific configuration stores, organizations could manage governed agentic configurations through repositories based on the ACM reference model. Such repositories could maintain complete configuration histories, provenance relationships, governance states, assurance evidence, and lifecycle information independently of the execution technologies used to deploy the corresponding systems. This approach would naturally extend established Software Configuration Management principles to heterogeneous agentic environments.

The emergence of a common governance ecosystem would also facilitate interoperability between independent software components. Framework adapters, governance engines, observability platforms, evaluation systems, compliance tools, and deployment infrastructures could exchange governance information through stable configuration abstractions rather than through framework-specific integration mechanisms. Such interoperability would reduce coupling between operational tools while improving the traceability and consistency of governance information across the software lifecycle.

Another promising direction concerns the definition of open interfaces supporting governance services. Standardized APIs for querying governed configurations, retrieving provenance information, validating lifecycle consistency, or evaluating governance states could enable independent tools to interact with ACM-compliant repositories without requiring knowledge of the underlying orchestration technologies. This service-oriented approach would encourage modular ecosystem development while preserving the framework independence of the reference model.

Finally, the emergence of an ecosystem around ACM would enable richer governance workflows spanning the complete lifecycle of agentic systems. Configuration management, execution monitoring, evaluation, assurance, auditing, compliance verification, and operational analytics could progressively become interconnected through a common governance representation while remaining implemented by specialized tools. In this architecture, ACM would serve as the semantic backbone connecting otherwise independent governance activities rather than replacing the specialized capabilities of existing AgentOps platforms.

Although such an ecosystem remains prospective, its development represents a natural continuation of the framework-independent philosophy introduced by ACM. By providing stable governance abstractions that are independent of execution technologies, the proposed reference model establishes the conceptual foundations upon which interoperable governance services, reusable software components, and collaborative governance infrastructures may progressively emerge.

\subsection{Standardization Perspectives}

The increasing diversity of agentic frameworks, operational platforms, and governance tools raises an important long-term question for the software engineering community: whether a common configuration governance model could facilitate interoperability across heterogeneous agentic ecosystems. While the present work does not propose a standard, the ACM reference model has been deliberately designed around framework-independent abstractions that could support future standardization initiatives.

Historically, Software Configuration Management evolved from collections of tool-specific practices toward broadly accepted conceptual models describing configuration items, revisions, baselines, change management, and traceability. These common abstractions enabled interoperability between heterogeneous tools while allowing implementation technologies to evolve independently. A comparable evolution may emerge for agentic systems as orchestration frameworks, AgentOps platforms, observability infrastructures, and governance services continue to diversify.

From this perspective, ACM may be viewed as a candidate open specification for configuration governance rather than as a framework-specific solution. Its explicit separation between governance semantics and execution semantics provides a stable conceptual layer upon which independent implementations can interoperate while preserving their own execution models and operational characteristics. Such a separation could facilitate the exchange of governed configurations, provenance information, lifecycle states, and governance metadata across heterogeneous software ecosystems.

Future standardization efforts would naturally require substantially broader validation than that presented in this work. The governance abstractions proposed by ACM would need to be evaluated across a significantly larger variety of orchestration frameworks, AgentOps platforms, industrial use cases, and application domains. In addition, long-term industrial experience would be necessary to assess whether the proposed concepts remain sufficiently stable as agentic technologies continue to evolve.

Another important prerequisite concerns community participation. Sustainable reference models and technical standards generally emerge through collaboration among researchers, software vendors, industrial users, and standardization bodies rather than from individual research contributions. Future evolution of ACM would therefore benefit from open discussions regarding governance terminology, semantic interoperability, extension mechanisms, compliance profiles, and implementation guidelines developed collectively by the broader agentic software engineering community.

Should such convergence progressively emerge, ACM could contribute to the establishment of common governance interfaces, exchange formats, conformance criteria, and interoperability profiles for governed agentic systems. These elements would facilitate independent implementations while preserving the flexibility required by heterogeneous execution technologies and rapidly evolving AI ecosystems.

Ultimately, the objective of ACM is not to prescribe a universal governance solution but to provide a formally defined reference model upon which future collaborative efforts may build. Whether this model eventually contributes to an open specification, an interoperability framework, or a broader standardization initiative will depend on its adoption, validation, and progressive refinement by the research and industrial communities. The present work should therefore be regarded as an initial contribution toward a shared conceptual foundation for configuration governance in agentic systems rather than as the endpoint of that process.

\section{Conclusion}

The rapid emergence of agentic AI has profoundly transformed the way intelligent software systems are designed and deployed. While orchestration frameworks have considerably simplified the construction of increasingly sophisticated agentic workflows, configuration governance has remained largely framework-specific, limiting reproducibility, traceability, interoperability, and long-term maintainability. This fragmentation motivated the development of the Agentic Configuration Management (ACM) reference model presented in this paper, whose objective is not to replace existing orchestration technologies but to provide a unified governance abstraction independent of execution environments.

This work introduced a framework-independent reference model for configuration governance in agentic systems. The proposed model defines a canonical representation of governed configuration items, immutable revisions, lifecycle semantics, provenance relationships, dependency structures, governance propagation mechanisms, and runtime representations. Unlike existing orchestration frameworks, ACM explicitly separates execution semantics from configuration governance, allowing heterogeneous agentic technologies to be represented within a common governance model while preserving their own operational characteristics.

Beyond the conceptual model, the paper proposed a formal mathematical semantics describing lifecycle management, governance state propagation, runtime evolution, and dynamic agent creation. This formalization establishes a precise semantic foundation that reduces ambiguity in the interpretation of governance operations and provides a basis for future theoretical analyses and formal verification. A reference implementation further demonstrated the practical feasibility of these concepts through deterministic governance mechanisms and immutable configuration management.

The experimental evaluation conducted using LangGraph and CrewAI demonstrated that heterogeneous orchestration models can be projected onto a common governance representation while preserving their essential governance semantics. Although the validation remains intentionally limited to two representative frameworks, the results provide initial evidence that configuration governance can be treated independently from orchestration technologies. This distinction constitutes one of the central contributions of the proposed reference model and suggests that governance abstractions may remain stable despite the rapid evolution of execution frameworks.

More broadly, this work argues that configuration governance should become a first-class concern of agentic software engineering. As Software Configuration Management established common principles for traditional software systems, agentic software engineering may similarly benefit from framework-independent governance abstractions capable of supporting reproducibility, provenance, auditability, interoperability, and long-term lifecycle management. Rather than introducing another orchestration framework, ACM proposes a semantic foundation upon which governance services, operational tooling, and heterogeneous AgentOps ecosystems may progressively interoperate.

Naturally, the ACM reference model represents only an initial step toward this objective. Broader validation across additional frameworks, richer industrial case studies, formal verification of the proposed semantics, and community-driven refinement will be necessary to assess the long-term applicability of the model and to determine whether its abstractions are sufficiently robust to support wider adoption. The future evolution of configuration governance for agentic systems will ultimately depend on collaborative contributions from both the research community and industrial practitioners.

In conclusion, this paper presents, to the best of our knowledge, one of the first attempts to define a formal, framework-independent reference model dedicated specifically to configuration governance for agentic systems. By establishing a clear separation between governance and execution semantics, introducing mathematically defined governance abstractions, and demonstrating their practical applicability through a working implementation and experimental validation, this work lays the foundations for a more interoperable and rigorous approach to the engineering of agentic systems. We hope that the ACM reference model presented in this paper will stimulate further research and contribute to the progressive emergence of shared governance principles for the next generation of agentic software engineering.






\section{Annexe}
\label{sec:annexe}
\begin{longtable}{p{0.9cm}p{2cm}p{6.2cm}p{2.3cm}p{1.6cm}}
\caption{Complete normative scenario catalogue.}
\label{tab:appendix_scenarios}\\

\toprule
ID &
Priority &
Normative objective &
Requirements &
Result \\
\midrule
\endfirsthead

\toprule
ID &
Priority &
Normative objective &
Requirements &
Result \\
\midrule
\endhead

S01 &
P0 &
Revision identity, immutable revisions and digest reproducibility &
I1--I4 &
Pass \\

S02 &
P0 &
Missing dependency detection and baseline consistency &
S02 &
Pass \\

S03 &
P0 &
Immutable revision history and append-only evolution &
S03 &
Pass \\

S04 &
P0 &
Baseline reproducibility from immutable revisions &
S04 &
Pass \\

S05 &
P0 &
Propagation of effective quality through dependencies &
S05 &
Pass \\

S06 &
P0 &
Eligibility propagation across dependent configuration items &
S06 &
Pass \\

S07 &
P0 &
Assurance coverage validation and dependency composition &
S07 &
Pass \\

S08 &
P1 &
Mixed quality propagation within heterogeneous dependency graphs &
S08 &
Pass \\

S09 &
P1 &
Propagation determinism under different traversal orders &
S09 &
Pass \\

S10 &
P1 &
Circular dependency detection and convergence guarantees &
S10 &
Pass \\

S11 &
P1 &
Fixed-point convergence of governance propagation &
S11 &
Pass \\

S12 &
P0 &
Runtime replay from immutable event log &
S12 &
Pass \\

S13 &
P0 &
Lifecycle transition validation during runtime replay &
S13 &
Pass \\

S14 &
P1 &
Declared runtime extension detection &
S14 &
Pass with deviation \\

S15 &
P1 &
Untraceable runtime instance detection &
S15 &
Pass with deviation \\

S16 &
P1 &
Configuration drift classification &
S16 &
Pass with deviation \\

S17 &
P0 &
Deterministic reconstruction of runtime state &
S17 &
Pass \\

S18 &
P1 &
Dynamic agent declaration and governance integration &
S18 &
Pass \\

S19 &
P1 &
Permission ceiling enforcement &
S19 &
Pass \\

S20 &
P1 &
Behaviour override validation &
S20 &
Pass \\

S21 &
P1 &
Promotion of dynamic agents into managed configuration &
S21 &
Pass \\

S22 &
P2 &
Projection consistency across supported frameworks &
S22 &
Pass \\

S23 &
P2 &
Framework-independent governance computation &
S23 &
Pass \\

S24 &
P2 &
Ordering invariance of normalized configuration &
S24 &
Pass \\

S25 &
P2 &
Robustness against incomplete framework metadata &
S25 &
Pass \\

S26 &
P2 &
Deterministic normalization of equivalent configurations &
S26 &
Pass \\

S27 &
P2 &
Local scalability on representative configuration graphs &
S27 &
Pass \\

\bottomrule
\end{longtable}

\begin{table}[t]
\centering
\caption{Provenance categories for normalized ACM information.}
\label{tab:information_provenance}
\begin{tabularx}{\linewidth}{lXX}
\toprule
Provenance category &
Definition &
Example \\
\midrule

Extracted &
Information obtained directly from an introspectable native framework object. &
LangGraph node or CrewAI agent \\

Declared by adapter &
Information explicitly supplied by the adapter when the native framework does not expose an equivalent static representation. &
Reconstructed CrewAI relationship \\

Derived &
Information deterministically computed from normalized ACM configuration or runtime data. &
Effective assurance or drift classification \\

Runtime-generated &
Information produced by an observed runtime event and appended to the event log. &
Agent creation or lifecycle transition \\

Unresolved &
Information whose value or semantics cannot be established from the available native objects or events. &
Opaque routing-function semantics \\

\bottomrule
\end{tabularx}
\end{table}
This provenance classification does not constitute a confidence or trust
assessment. It records how each item entered the ACM representation and enables
auditors to distinguish directly observed information from adapter declarations,
algorithmically derived values, runtime observations and unresolved semantics.
Assessing the reliability of these sources would require an additional
validation model and is outside the scope of the present evaluation.

\end{document}
